\section{Conventions}
\label{ch:conventions}

This chapter collects the conventions on which every other chapter depends: the state space and the fields, the normalization of the sphere correlators, the pairings on one and several copies of the state space with their exchange signs, and the reality conditions on the string field and on the auxiliary fields.
It expands the two convention appendices of the paper, which are reproduced first and in full.
They are followed by four sections of an older companion note on the pairing and the grading, written in an earlier notation that the translation box before them relates to the paper's.
The chapter ends with three sections that exist only here: the conversion between the two normalizations of the sphere correlator met in this project, the orientation of the off-shell measure, and the derivation of the symmetry type of the strands of a multi-string tensor, which the paper states without proof.

\paragraph*{Reading note.}
\status{Verified} The first section is the convention appendix of the paper.
The normalization of the sphere correlator it fixes is the one of~\cite{Erbin:2021:StringFieldTheory}; how numbers computed in the other common normalization are converted is the subject of \cref{sec:conv:two-normalizations}.\record{041-convention-reconciliation (R43)}
\status{mixed} The reality conditions and the equivariance signs in the star convention, at the end of the section, have different statuses.\record{075-star-conjugation-convention (R81)}
The definition of the star conjugation, its relation to the conjugation used in the older notes, the table of reality conditions and the equivariance of the pairing and of the kinetic term are proven by hand.
The reality of the vertices with three or more arguments, and with it the interaction part of the condition on the gauge parameter, rests on the standing assumption on the vertices stated in the section, and is argued rather than checked; the condition on the multipliers $\Theta_k$ is checked in a finite model at $k = 2, 3$ and argued in general.
The matter sector is argued from the rule on the modes, the ghost zero-mode sector is checked explicitly.

\subsection{Conventions and normalization}
\label{p:app:conventions}

This appendix collects the conventions of the paper, following~\cite{Erbin:2021:StringFieldTheory}.


\subparagraph{State space.}

The zero modes of the ghosts and of the Virasoro generators are combined into
\begin{equation}
  \label{p:eq:conv:pm-operators}
  L_0^\pm
    := L_0 \pm \bar L_0,
  \qquad
  b_0^\pm
    := b_0 \pm \bar b_0,
  \qquad
  c_0^\pm
    := \frac{1}{2} \, (c_0 \pm \bar c_0).
\end{equation}
$\Hs$ denotes the space of states of the matter--ghost CFT, at every ghost number, that satisfy the level-matching conditions
\begin{equation}
  \label{p:eq:conv:level-matching}
  b_0^- \ket{A}
    = 0,
  \qquad
  L_0^- \ket{A}
    = 0.
\end{equation}
We include all ghost numbers so that the string products below map into $\Hs$.

Ghost numbers are physical: the $\group{SL}(2, \C)$ vacuum has ghost number zero, $c$ counts $+1$ and $b$ counts $-1$.
The tachyon $c_1 \bar c_1 \ket{0}$ then has ghost number two.
The Grassmann parity of a state is its ghost number modulo two:
\begin{equation}
  \label{p:eq:conv:parity}
  \abs{A}
    := \Ngh(A) \bmod 2.
\end{equation}

The $n$-string space $\Hn{n}$ is the space of tensors in $\Hs^{\otimes n}$ whose components are graded-symmetric under permutations of the strands.
Given a basis $\{ \phi_r \}$ of $\mc H$, we can write a state $\Sigma_n \in \mc H^{\wedge n}$ as
\begin{equation}
  \label{p:eq:conv:exp-sigma-n}
  \Sigma_n
    = \sum_{\{ r_i \}} \sigma_{r_1 \cdots r_n} \phi_{r_1} \cdots \phi_{r_n}.
\end{equation}

The parity in which the symmetry is graded depends on the role of the strands.
For the arguments of the vertices and products below, the fields $\Psi^n$ and $\Sigma_n$ and the gauge parameters, the grading is given by the parity $\abs{\phi_{r_i}}$ of the basis state in the strand $i$ of~\eqref{p:eq:conv:exp-sigma-n}, and $\Hn{n}$ is the symmetrized tensor product of $n$ copies of $\mc H$~\cite{Erler:2014:NSNSSectorClosed,Erbin:2020:ClassicalAlgebraicStructures}.
Even strands then have symmetric components and odd strands antisymmetric ones, so that a product is antisymmetric in two gauge parameters.

For the strands that the $n$-string pairing~\eqref{p:eq:conv:pairing-n} sets against a field, those of the outputs of the products and of the multipliers $\Theta_k$, it is the shifted parity $\abs{\phi_{r_i}} + 1$.
This is the dual of the former space under the pairing: if we were also defining the graded-symmetry with respect to the strand parity, the outputs at ghost number three would have antisymmetric components and would pair to zero with $\Sigma_n$.
For this reason, we use instead the shifted parity $\abs{\phi_{r_i}} + 1$.
In this way, the fields have symmetric components, as the pairing requires.
For a tensor with one odd and one even strand, the two parities give the same sign and we need not specify the role.
We write $\Psi^n := \Psi^{\otimes n} \in \Hn{n}$.


\subparagraph{Fields of the paper.}

We list below the fields of the paper with their space, ghost number and Grassmann parity, and the section where each is introduced.
All of them have commuting coefficients, so that the parity of a field is that of its states.
\begin{itemize}
  \item $\Psi \in \Hs$, the classical string field: $\Ngh(\Psi) = 2$, even (\cref{p:sec:hs:action}).

  \item $\Lambda \in \Hs$, the classical gauge parameter: $\Ngh(\Lambda) = 1$, odd (\cref{p:sec:hs:action}).

  \item $\Sigma_k \in \Hn{k} / \ker$, the auxiliary field that replaces the tensor power $\Psi^{\otimes k}$, on the quotient by the kernel of its quadratic operator, $\sprod{2}{2}$ for $k = 2$ and $\wprod{3}{3}$ for $k = 3$: $\Ngh(\Sigma_k) = 2 k$, even (\cref{p:sec:hs:quartic,p:sec:hs:sextic}).
    In the multiplier form of \cref{p:sec:all-orders}, the fields $\Sigma_n \in \Hn{n}$ carry the same letter but live on $\Hn{n}$ itself and are not these auxiliary fields.

  \item $\Theta_k \in \Hn{k}$, the multiplier that enforces $\Sigma_k = \Sigma_{k-1} \otimes \Psi$ in the multiplier form of the action: $\Ngh(\Theta_k) = 3 k$, parity $k \bmod 2$, its strands read in the shifted parity (\cref{p:sec:all-orders}).

  \item $\Upsilon_2 \in \Hn{2} / \ker \sprod{2}{2}$, the partner field: $\Ngh(\Upsilon_2) = 4$, even (\cref{p:app:wrong-sign}).
\end{itemize}
Three letters should not be confused: $d$ is the spacetime dimension, $D$ the ghost dilaton of \cref{p:sec:solutions}, and $\hat d$ the formal dimension of \cref{p:sec:worldsheet}.


\subparagraph{Sphere normalization.}

We normalize the correlation functions on the sphere, with the momentum-conservation factor $(2\pi)^d \delta^{(d)}(\sum k)$ stripped off, as:
\begin{equation}
  \label{p:eq:conv:sphere}
  \braket{c_{-1} \bar c_{-1} c_0^- c_0^+ c_1 \bar c_1}
    = - \frac{1}{2},
  \qquad
  \braket{c_{-1} \bar c_{-1} c_0 \bar c_0 c_1 \bar c_1}
    = - 1.
\end{equation}
The value is real, with no factor $\I$, as the reality of the action under~\eqref{p:eq:conv:reality-fields} requires.
A sphere correlator is nonzero only when the ghost numbers of its insertions add up to six, the ghost-number anomaly.
All the ghost-number formulas below follow from this selection rule.


\subparagraph{Pairing.}

The pairing of two states is the BPZ inner product with one insertion of $c_0^-$:
\begin{equation}
  \label{p:eq:conv:pairing}
  \pair{A}{B}
    := \braket{A | c_0^- | B},
\end{equation}
where $\bra{A}$ is the BPZ conjugate of $\ket{A}$.
We write the BPZ conjugate as $\transp{A}$.
The pairing is nonzero only when $\Ngh(A) + \Ngh(B) = 5$, since the insertion of $c_0^-$ supplies the sixth unit.
Hence it pairs states of opposite parity, $\abs{A} + \abs{B} = 1$.

BPZ conjugation acts on the modes as
\begin{equation}
\label{p:eq:conv:bpz-modes}
\transp{c_n}
  = - c_{-n},
\qquad
\transp{b_n}
  = b_{-n},
\qquad
\transp{(c_0^\pm)}
  = - c_0^\pm,
\end{equation}
and in the same way for $\bar c_n$ and $\bar b_n$.
These are the rules for the inversion $I^+(z) = 1/z$, the standard choice for the closed string~\cite[footnote p.~13]{Zwiebach:1993:ClosedStringField}.
On a monomial of creation modes acting on the vacuum, BPZ conjugation conjugates each mode in place, without a reordering sign:
\begin{equation}
  \label{p:eq:conv:bpz-word}
  \transp{(O_1 \cdots O_k \ket{0})}
    = \bra{0} \transp{O_1} \cdots \transp{O_k}.
\end{equation}

With these rules and the normalization~\eqref{p:eq:conv:sphere}, the pairing is graded-symmetric in the shifted parity:
\begin{equation}
\label{p:eq:conv:pairing-symmetry}
\pair{B}{A}
  = (-1)^{(1 + \abs{A}) (1 + \abs{B})} \, \pair{A}{B}.
\end{equation}
Since the pairing is nonzero only for $\abs{A} + \abs{B} = 1$, the sign is $+1$ on its support: the pairing is symmetric.
This is the rearrangement rule of~\cite{Erbin:2021:StringFieldTheory} restricted to the support.
The sign is checked on explicit states in \cref{p:app:reality}, where the reality conditions are also derived.


\subparagraph{Multi-string pairing.}

On $n$-string states the pairing is the tensor power of~\eqref{p:eq:conv:pairing}, one $c_0^-$ per strand, with the Koszul sign of a tensor product of bilinear forms,
\begin{equation}
  \label{p:eq:conv:pairing-n}
  \pair[n]{A_1 \otimes \cdots \otimes A_n}{B_1 \otimes \cdots \otimes B_n}
    := (-1)^{\sum_{i < j} \abs{B_i} \abs{A_j}} \,
      \pair{A_1}{B_1} \cdots \pair{A_n}{B_n},
\end{equation}
extended bilinearly.
For $n = 1$ it is the pairing~\eqref{p:eq:conv:pairing}, written without index.
When the strands of the first tensor are even, as for $\Psi^n$ and $\Sigma_n$, the sign is $+1$ and the pairing is the plain product of the pairings of the strands.
Each $c_0^-$ thus sits inside the pairing of its own strand, and we read the products in this grouping.
The multi-string pairing is nonzero only when the ghost numbers of the two tensors add up to $5n$.

For every $n$, it is graded-symmetric in the total parity:
\begin{equation}
\label{p:eq:conv:pairing-n-symmetry}
\pair[n]{Y}{X}
  = (-1)^{\abs{X} \abs{Y}} \, \pair[n]{X}{Y},
\qquad
X, Y \in \Hn{n}.
\end{equation}
This follows from combining the Koszul signs with~\eqref{p:eq:conv:pairing-symmetry} strand by strand.
Whenever one of the two tensors is even, as for every pair met in this paper, the sign is $+1$ and the pairing is symmetric.


\subparagraph{String vertices.}

The vertices $\vertex{n} \colon \Hn{n} \to \C$ are the genus-zero string vertices, graded-symmetric multilinear forms.
Their arguments, and those of the products, are graded-symmetric in the parity of the strands for every parity of the arguments.
Hence the product of two gauge parameters in \cref{p:sec:gauge} is antisymmetric in the two.
The outputs of the multi-string products are read in the shifted parity.
The string products $\ell_n \colon \Hn{n} \to \Hs$ are the maps they define through the pairing:
\begin{equation}
  \label{p:eq:conv:vertex}
  \vertex{n}(A_1, \ldots, A_n)
    =: \pair{A_1}{\ell_{n-1}(A_2, \ldots, A_n)},
  \qquad
  \ell_1
    := Q_B.
\end{equation}
For $n = 2$, we have in particular, $\vertex{2}(A_1, A_2) := \pair{A_1}{Q_B A_2}$.
The BRST operator is graded self-adjoint for the pairing:
\begin{equation}
  \label{p:eq:conv:QB-self-adjoint}
  \pair{Q_B A}{B}
    = (-1)^{\abs{A}} \, \pair{A}{Q_B B}.
\end{equation}
This property is the symmetry of $\vertex{2}$.

A vertex extended by multilinearity to tensor arguments is written with the tensors in its slots: using~\eqref{p:eq:conv:exp-sigma-n},
\begin{equation}
  \vertex{n}(\Sigma_k, \cdots)
    := \sum_{\{ r_i \}} \sigma_{r_1 \cdots r_k} \, \vertex{n}(\phi_{r_1}, \ldots, \phi_{r_k}, \cdots).
\end{equation}

The selection rule fixes the ghost number carried by every map.
For $\ell_n$ and $\sprod{m}{n}$ it is the ghost number of the output minus the sum of those of the inputs.
For $\vertex{n}$ it is the ghost number of its antighost insertions, so that $\vertex{n}(A_1, \ldots, A_n)$ is nonzero only when $\sum_i \Ngh(A_i) = 2n$:
\begin{equation}
  \label{p:eq:conv:ghost-numbers}
  \Ngh(\vertex{n})
    = 6 - 2n,
  \qquad
  \Ngh(\ell_n)
    = 3 - 2n,
  \qquad
  \Ngh(\sprod{m}{n})
    = 3n - 2m.
\end{equation}
The last entry concerns the products with several outputs, $\sprod{m}{n} \colon \Hn{m} \to \Hn{n}$, whose output strands have ghost number three.
They are defined by splitting the slots of a vertex into $n$ outputs and $m$ inputs and reading the result through the $n$-string pairing,
\begin{equation}
  \label{p:eq:conv:sprod}
  \vertex{m+n}(A_1, \ldots, A_n, B_1, \ldots, B_m)
    =: \pair[n]{A_1 \otimes \cdots \otimes A_n}{\sprod{m}{n}(B_1, \ldots, B_m)},
\end{equation}
so that $\sprod{n}{1} = \ell_n$ and, for instance:
\begin{equation*}
  \vertex{4}(A_1, A_2, A_3, A_4)
    = \pair[2]{A_1 \otimes A_2}{\sprod{2}{2}(A_3, A_4)}.
\end{equation*}
Each output raises the ghost number by three and each input lowers it by two, so that $\Ngh(\sprod{2}{2}) = 2$, $\Ngh(\sprod{3}{2}) = 0$ and $\Ngh(\sprod{3}{3}) = 3$.

Since $\vertex{m+n}$ is graded-symmetric, reading the same vertex with inputs and outputs exchanged gives the transposition identity
\begin{equation}
  \label{p:eq:conv:transpose}
  \pair[n]{X}{\sprod{m}{n}(Y)}
    = \pair[m]{\sprod{n}{m}(X)}{Y},
  \qquad
  X \in \Hn{n},
  \quad
  Y \in \Hn{m},
\end{equation}
for even arguments, the case of $\Psi$ and the auxiliary fields (odd arguments bring a sign).
$\sprod{n}{m}$ is thus the transpose of $\sprod{m}{n}$ with respect to the multi-string pairing.
We write compositions with $\circ$.

\subparagraph{Normalization of the action.}

With the vertices normalized as in~\eqref{p:eq:conv:vertex}, the action of closed string field theory is:
\begin{equation}
  \label{p:eq:conv:action}
  S
    = \frac{1}{2} \, \pair{\Psi}{Q_B \Psi}
      + \sum_{n \ge 3} \frac{\gs^{n-2}}{n!} \, \vertex{n}(\Psi^n)
    = \sum_{n \ge 1} \frac{\gs^{n-1}}{(n+1)!} \, \pair{\Psi}{\ell_n(\Psi^n)},
\end{equation}
with $\gs$ the string coupling.
The equation of motion and the gauge transformation read:
\begin{equation}
  \label{p:eq:conv:gauge}
  Q_B \Psi
    + \sum_{n \ge 2} \frac{\gs^{n-1}}{n!} \, \ell_n(\Psi^n)
    = 0,
  \qquad
  \delta_\Lambda \Psi
    = Q_B \Lambda
      + \sum_{n \ge 2} \frac{\gs^{n-1}}{(n-1)!} \, \ell_n(\Psi^{n-1}, \Lambda),
\end{equation}
with $\Lambda \in \Hs$ of ghost number one.

The Euclidean path integral is weighted by $\e^{S}$, that is, $S = -S_E$ with $S_E$ the Euclidean action.
  For the tachyon component $t(k)$ of the string field, $\Psi \supset t(k) \, c_1 \bar c_1 \ket{k}$, with $m_T^2 = -4/\alpha'$, the quadratic term of~\eqref{p:eq:conv:action} is:
  \begin{equation}
  \label{p:eq:conv:tachyon-quadratic}
  S
    \supset - \frac{\alpha'}{8} \int \frac{\dd^d k}{(2\pi)^d} \, (k^2 + m_T^2) \, t(k) \, t(-k).
\end{equation}
  It decreases at large Euclidean momentum, as a damped weight requires.
  This fixes the sign of~\eqref{p:eq:conv:sphere} against the sign of the weight.


\subparagraph{Reality.}


We impose reality with the star conjugation of~\cite{Erbin:2021:StringFieldTheory}, the composition of the BPZ and Hermitian conjugations:
\begin{equation}
  \label{p:eq:conv:star}
  A^\star
    := \adj{(\transp{A})}
    = \transp{(\adj{A})}.
\end{equation}
The BPZ bra of $A^\star$ is thus the Hermitian bra of $A$, and conversely:
  \begin{equation}
  \label{p:eq:rea:star-bras}
  \bra{A^\star}
    = \adj{A},
  \qquad
  \adj{(A^\star)}
    = \bra{A},
\end{equation}
  with $\bra{A} = \transp{A}$.
  The two conjugations commute on every mode, since each maps $c_n$, $b_n$ and $\alpha_n$ to a real multiple of the mode of opposite index.
  On the modes, Hermitian conjugation is $\adj{(c_n)} = c_{-n}$ and $\adj{(b_n)} = b_{-n}$, and it reverses the order of a product of modes.
  On operators the star is the same composition, and it distributes over the action on a state with a Koszul sign:
  \begin{equation}
  \label{p:eq:rea:star-operator}
  O^\star
    := \transp{(\adj{O})},
  \qquad
  (O A)^\star
    = (-1)^{\abs{O} \abs{A}} \, O^\star A^\star.
\end{equation}
  The sign appears because $O^\star$ is written to the left of $A^\star$ while Hermitian conjugation reverses the order.
Star conjugation is antilinear and squares to one on every state, even or odd: on a monomial of modes acting on a momentum state it returns the same monomial, up to a sign, on the opposite momentum.
On a tensor it acts strand by strand, without a reordering sign.
On the string field, the reality condition is that of~\cite[eq.~(3.6)]{Zwiebach:1993:ClosedStringField}:
\begin{equation}
  (\ket{\Psi})^\dagger
    = - \bra{\Psi}.
\end{equation}

The pairing is equivariant up to a minus sign,
\begin{equation}
  \label{p:eq:conv:pairing-equivariance}
  \pair{A^\star}{B^\star}
    = - \, \pair{A}{B}^*,
\end{equation}
for all $A$ and $B$.
Complex conjugation of a sphere correlator and Hermitian conjugation of its insertions define two antilinear functionals on the monomials of modes, which vanish on the same monomials.
  Since the vacuum functional is unique up to a factor, there is a constant $\eta$ such that $\braket{W}^* = \eta \, \braket{\adj{W}}$ for every monomial $W$.
  On $W_0 = c_{-1} \bar c_{-1} c_0 \bar c_0 c_1 \bar c_1$, Hermitian conjugation swaps each pair $(c_n, \bar c_n)$, which amounts to three transpositions of anticommuting modes, so that $\adj{W_0} = - W_0$.
  With $\braket{W_0} = - 1$ from~\eqref{p:eq:conv:sphere}, this gives $\eta = - 1$: the vacuum functional is anti-Hermitian, as in~\cite[eq.~(2.37)]{Zwiebach:1993:ClosedStringField}.
  Writing $A = U \ket{0}$ and $B^\star = V \ket{0}$ with monomials of modes $U$ and $V$, the relations~\eqref{p:eq:rea:star-bras} and $\adj{(c_0^-)} = c_0^-$ give $\pair{A^\star}{B^\star} = \eta \, \pair{B}{A}^*$, which is~\eqref{p:eq:conv:pairing-equivariance} by the symmetry~\eqref{p:eq:conv:pairing-symmetry}.
The BRST operator is Hermitian and odd under BPZ conjugation, so that $Q_B^\star = - Q_B$.
  With~\eqref{p:eq:rea:star-operator} and $\abs{Q_B} = 1$, it commutes with star conjugation up to a sign fixed by the parity of its argument:
\begin{equation}
  \label{p:eq:conv:QB-star}
  (Q_B A)^\star
    = (-1)^{\abs{A} + 1} \, Q_B A^\star.
\end{equation}
The vertices are real in the sense that
\begin{equation}
  \label{p:eq:conv:vertex-reality}
  \vertex{n}(A_1, \ldots, A_n)^*
    = (-1)^{n + \sum_{i < j} \abs{A_i} \abs{A_j}} \, \vertex{n}(A_1^\star, \ldots, A_n^\star).
\end{equation}
On even arguments, the sign is $(-1)^n$.
  This property makes the action~\eqref{p:eq:conv:action} real.
For $n = 2$ it follows from~\eqref{p:eq:conv:pairing-equivariance} and~\eqref{p:eq:conv:QB-star}, since $\vertex{2}(A_1, A_2) = \pair{A_1}{Q_B A_2}$.
For $n \ge 3$, we take it as a standing assumption on the vertices.
It follows from the invariance of the vertex regions and measures under complex conjugation of the moduli together with the hermiticity conditions of~\cite[sec.~5.1 and eq.~(7.98)]{Zwiebach:1993:ClosedStringField}.
For the $\group{SL}(2, \C)$ vertices of~\cite{Erbin:2026:SL2CVerticesClosed} it holds for the quartic vertex, and we assume it for every $n$.

Since star conjugation acts strand by strand and preserves the parity of every strand, the multi-string pairing is equivariant as well, $\pair[n]{X^\star}{Y^\star} = (-1)^n \, \pair[n]{X}{Y}^*$.
  The reality conditions on the fields of the paper are:
\begin{equation}
  \label{p:eq:conv:reality-fields}
  \Psi^\star
    = - \Psi,
  \qquad
  \Sigma_k^\star
    = (-1)^k \, \Sigma_k,
  \qquad
  \Lambda^\star
    = - \Lambda.
\end{equation}
The condition on $\Psi$ is the reality condition of the theory, and the condition on $\Lambda$ is the one under which every gauge transformation~\eqref{p:eq:conv:gauge} preserves it~\cite[eq.~(3.18)]{Zwiebach:1993:ClosedStringField}, the free term by the star of $Q_B$ and the interaction terms by the reality of the vertices (\cref{p:app:reality}).
  The condition on $\Sigma_k$ follows from the reality of $\Psi$ and of the vertices once the auxiliary fields are introduced (\cref{p:app:reality}).


\subparagraph{Batalin--Vilkovisky conventions.}

In the Batalin--Vilkovisky formalism of \cref{p:sec:gauge}, every field $\phi^i$ is paired with an antifield $\phi^+_i$ of opposite Grassmann parity and target-space ghost number $\gh(\phi^+_i) = -\gh(\phi^i) - 1$.
The master action is even, of ghost number zero.
In the standard assignment of~\cite{Zwiebach:1993:ClosedStringField}, the antifield of a component of $\Psi$ at worldsheet ghost number $g$ sits at ghost number $5 - g$, so that the two pair under~\eqref{p:eq:conv:pairing}.
For a $2$-string field the pairing is~\eqref{p:eq:conv:pairing-n}, nonzero when the ghost numbers add up to ten, so the antifield of $\Sigma_2$ sits at ghost number six.
We take the antibracket in the convention of~\cite[ch.~15]{Henneaux:1994:QuantizationGaugeSystems},
\begin{equation}
  \label{p:eq:conv:antibracket}
  (F, G)
    := \frac{\pd^r F}{\pd \phi^i} \, \frac{\pd^l G}{\pd \phi^+_i}
      - \frac{\pd^r F}{\pd \phi^+_i} \, \frac{\pd^l G}{\pd \phi^i},
\end{equation}
with right and left derivatives as indicated.

On the string field, the odd symplectic form is the pairing twisted by the parity of its first argument:
\begin{equation}
  \label{p:eq:conv:odd-symplectic}
  \Omega(A, B)
    := (-1)^{\abs{A}} \, \pair{A}{B},
  \qquad
  \Omega(B, A)
    = - (-1)^{\abs{A} \abs{B}} \, \Omega(A, B).
\end{equation}
It is odd and graded-antisymmetric on every pair of states in its support, since the twist changes sign under the exchange of two arguments of opposite parity, whereas the pairing itself is symmetric.
Twisting by the parity of the second argument gives $-\Omega$.
Its overall sign relative to the antibracket~\eqref{p:eq:conv:antibracket} is a convention fixed where the master action is written.

On the auxiliary fields, the parity $\Ngh \bmod 2$ does not give an odd pairing: $\Sigma_2$ and the $2$-string states it pairs with, at ghost numbers four and six, are both even in that grading.
We take the antifield $\Sigma_2^+$ of \cref{p:sec:gauge} to be a $2$-string state paired with $\Sigma_2$ by the $2$-string pairing, hence of worldsheet ghost number six.
It is Grassmann odd and of target-space ghost number $-1$ by the definition of an antifield.
Its worldsheet ghost number is a label on which nothing depends.
The rule~\eqref{p:eq:conv:parity} applies to the string field, its gauge parameter and their antifields.
On the auxiliary pair the overall sign of $\Omega$ is absorbed in the sign of $\Sigma_2^+$ (\cref{p:sec:gauge}).


\subparagraph{Potential.}


The potential is the action at zero momentum with the sign reversed,
\begin{equation}
  \label{p:eq:conv:potential}
  V
    := - S \big\vert_{k = 0},
\end{equation}
so that the tachyon sits at a maximum, $V = -\tfrac{1}{2} \, t^2$ at quadratic order.
A vertex quoted on the potential side carries the sign opposite to that of $\vertex{n}$.
When we count the negative directions of a Hessian, we say on which side it is taken, since the numbers of positive and negative directions of $\Hess S$ are those of $\Hess V$ exchanged.


\paragraph*{Reading note.}
\status{Proven} The next section is the second convention appendix of the paper.
It derives the exchange sign of the one-string pairing on explicit states and its extension to the $k$-string pairing: in the physical parity, $\pair[k]{Y}{X} = (-1)^{\abs{X} \abs{Y}} \pair[k]{X}{Y}$ for every $k$, the one-string pairing being symmetric on its support.\record{055-omega-k-graded-symmetry (R58)}
\status{Verified} The case $k = 2$ is the statement that the pairing of the auxiliary field $\Sigma_2$ is graded symmetric, hence not an odd symplectic form, in the physical grading; it is checked on finite models.\record{002-omega2-construction (R7)}
\status{mixed} The section ends with the reality of the auxiliary fields.\record{028-reality-tower (R35)}
Given the reality of the vertices, the reality of the whole tower of auxiliary fields off shell, with its kinetic terms, couplings and sources, is proven; at the solution it is proven under the kernel and invertibility conditions on the quadratic operators stated in the chapters on the construction, one of which is open; for the partner field it is proven through $O(\gs^4)$ and argued beyond.
The reality condition on $\Sigma_k$ is forced by the sources only up to the directions that no source reaches.

\subsection{Exchange sign of the pairing and reality of the auxiliary fields}
\label{p:app:reality}


This appendix checks the exchange sign~\eqref{p:eq:conv:pairing-symmetry} of the pairing on explicit states, with its consequences for the multi-string pairings, and derives the reality of the vertices, of the products and of the auxiliary fields.
We use throughout the normalization~\eqref{p:eq:conv:sphere} and the BPZ rules~\eqref{p:eq:conv:bpz-modes}.


\subsubsection{Exchange sign of the pairing}

The pairing $\pair{A}{B} = \braket{A | c_0^- | B}$ differs from the plain BPZ product $\braket{A | B}$, which is graded-symmetric in the parity of its arguments, by the odd insertion $c_0^-$.
Its exchange sign comes from moving the insertion from one argument to the other.
Moving an operator $O$ from the ket to the bra gives the sign of passing $O$ through $A$:
\begin{equation}
\braket{A | O B}
  = (-1)^{\abs{O} \abs{A}} \, \braket{\transp{O} A | B}.
\end{equation}
Applied to $O = c_0^-$, with $\transp{(c_0^-)} = - c_0^-$ from~\eqref{p:eq:conv:bpz-modes}, we get:
\begin{equation}
  \label{p:eq:rea:exchange-general}
  \begin{aligned}
    \pair{A}{B}
      &
      = \braket{A | c_0^- B}
      = (-1)^{\abs{A} (1 + \abs{B})} \, \braket{c_0^- B | A}
    \\
      &
      = - (-1)^{\abs{A} (1 + \abs{B})} \, (-1)^{\abs{B}} \, \pair{B}{A}
      = (-1)^{(1 + \abs{A}) (1 + \abs{B})} \, \pair{B}{A}.
  \end{aligned}
\end{equation}
The first sign comes from exchanging $A$ with the state $c_0^- B$, of parity $\abs{B} + 1$, in the plain product.
In the second step, we move $c_0^-$ into the bra, which costs the sign $(-1)^{\abs{B}}$ of passing the insertion through $B$ and the minus sign of its BPZ conjugate.
This is~\eqref{p:eq:conv:pairing-symmetry}.
The in-place rule~\eqref{p:eq:conv:bpz-word} brings no further sign: it is the action of the graded anti-involution on a set of creation operators, and a rule reversing the order without a graded sign would be inconsistent beyond the zero modes.

We check~\eqref{p:eq:rea:exchange-general} on the tachyon state $\ket{A} = c_1 \bar c_1 \ket{0}$ and its partner at ghost number three $\ket{B} = c_0^+ c_1 \bar c_1 \ket{0}$.
Their bras follow from~\eqref{p:eq:conv:bpz-modes} and~\eqref{p:eq:conv:bpz-word}, $\bra{A} = \bra{0} c_{-1} \bar c_{-1}$ and $\bra{B} = - \bra{0} c_0^+ c_{-1} \bar c_{-1}$.
  With the normalization~\eqref{p:eq:conv:sphere}, we get:
\begin{equation}
  \label{p:eq:rea:check-values}
  \begin{aligned}
    \pair{A}{B}
      &
      = \braket{c_{-1} \bar c_{-1} c_0^- c_0^+ c_1 \bar c_1}
      = - \frac{1}{2},
    \\
    \pair{B}{A}
      &
      = - \braket{c_0^+ c_{-1} \bar c_{-1} c_0^- c_1 \bar c_1}
      = \braket{c_{-1} \bar c_{-1} c_0^- c_0^+ c_1 \bar c_1}
      = - \frac{1}{2}.
  \end{aligned}
\end{equation}
In the second value, $c_0^+$ passes the even pair $c_{-1} \bar c_{-1}$ without a sign and then anticommutes with $c_0^-$.
The two values agree, as~\eqref{p:eq:conv:pairing-symmetry} requires for an even $A$ and an odd $B$.
  With $\transp{(c_0^\pm)} = + c_0^\pm$, the bra $\bra{B}$ would change sign and the two values would differ.
On the same pair, we have $A^\star = - A$, from the reversal of the two odd modes $c_1 \bar c_1$, and $B^\star = + B$, from the sign $- 1$ of $\transp{(c_0^+)}$ and the reversal sign $- 1$ of three odd modes.
  Hence $\pair{A^\star}{B^\star} = + \tfrac{1}{2} = - \pair{A}{B}^*$, as the equivariance~\eqref{p:eq:conv:pairing-equivariance} requires.

The $n$-string pairing~\eqref{p:eq:conv:pairing-n} inherits the symmetry strand by strand.
We write $X = x_1 \otimes \cdots \otimes x_n$ and $Y = y_1 \otimes \cdots \otimes y_n$.
Exchanging the two tensors exchanges the roles of the $x_i$ and the $y_j$ in the Koszul sign of the definition, and the arguments of every strand pairing:
\begin{equation}
  \label{p:eq:rea:exchange-n}
  \begin{aligned}
    \pair[n]{Y}{X}
      &
      = (-1)^{\sum_{i < j} \abs{x_i} \abs{y_j}} \, \prod_i \pair{y_i}{x_i}
      \\ &
      = (-1)^{\sum_{i < j} \abs{x_i} \abs{y_j} + \sum_i \abs{x_i} \abs{y_i} + \sum_{i < j} \abs{y_i} \abs{x_j}} \, \pair[n]{X}{Y}
      \\ &
      = (-1)^{\abs{X} \abs{Y}} \, \pair[n]{X}{Y}.
  \end{aligned}
\end{equation}
The three sums cover every pair $(i, j)$ exactly once, so the exponent is $\sum_{i, j} \abs{x_i} \abs{y_j} = \abs{X} \abs{Y}$.
This is~\eqref{p:eq:conv:pairing-n-symmetry}, on every pair of homogeneous tensors and by bilinearity on all of $\Hs^{\otimes n}$.

The same bookkeeping fixes the space that the auxiliary fields pair with.
Permuting the strands of both tensors by the same permutation $\pi$, each with the Koszul sign of the physical parity, multiplies the pairing on its support by $\sign(\pi)$.
Indeed, regrouped by pairs of strands, each inversion of $\pi$ exchanges two pairs of total parity $\abs{x_i} + \abs{y_i} = 1$ and costs one sign.
We take $X = \Sigma_n$, with even strands and symmetric components, and $Y$ with strands of ghost number three, all odd.
Then $\pair[n]{\Sigma_n}{Y} = \sign(\pi) \, \pair[n]{\Sigma_n}{\pi Y}$, and only the part of $Y$ with $\pi Y = \sign(\pi) \, Y$ pairs with $\Sigma_n$.
On all-odd strands, the Koszul action of $\pi$ is the plain permutation of the components times $\sign(\pi)$.
The part of $Y$ that pairs with $\Sigma_n$ is thus the part with plainly symmetric components, that is the graded-symmetric tensors in the shifted parity of \cref{p:app:conventions}.

The odd symplectic form~\eqref{p:eq:conv:odd-symplectic}, $\Omega(A, B) = (-1)^{\abs{A}} \, \pair{A}{B}$, is then graded-antisymmetric on every homogeneous pair by~\eqref{p:eq:rea:exchange-general}: the twist $(-1)^{\abs{A}}$ changes sign under the exchange of arguments of opposite parity.


\subparagraph{Reality of the vertices.}

For $n = 2$, the reality~\eqref{p:eq:conv:vertex-reality} follows from the equivariance~\eqref{p:eq:conv:pairing-equivariance} and the star of the BRST operator~\eqref{p:eq:conv:QB-star}:
\begin{equation}
  \label{p:eq:rea:V2-star}
  \begin{aligned}
    \vertex{2}(A^\star, B^\star)
      &
      = \pair{A^\star}{Q_B B^\star}
      = (-1)^{\abs{B} + 1} \, \pair{A^\star}{(Q_B B)^\star}
    \\
      &
      = (-1)^{\abs{B}} \, \pair{A}{Q_B B}^*
      = (-1)^{\abs{A} \abs{B}} \, \vertex{2}(A, B)^*,
  \end{aligned}
\end{equation}
since on the support $\abs{A} = \abs{B}$, so that $\abs{A} \abs{B} = \abs{B}$ modulo two.
Read term by term on the coefficients, with the reality condition~\cite[eq.~(3.6)]{Zwiebach:1993:ClosedStringField} on each of the $n$ fields, the reality of the vertices gives the factor $(-1)^n$ from the sign of that condition and the Koszul factor from reordering the coefficients.
We obtain the reality of the products by stripping the first argument of~\eqref{p:eq:conv:vertex} with~\eqref{p:eq:conv:pairing-equivariance}, since the pairing is non-degenerate on the level-matched space.
In $\vertex{m + 1}(A, B_1, \ldots, B_m) = \pair{A}{\ell_m(B_1, \ldots, B_m)}$, the output has parity $1 + \sum_i \abs{B_i}$.
On the support of the pairing, we thus have $\abs{A} = \sum_i \abs{B_i}$ modulo two, and the rule for arguments of any parity reads:
\begin{equation}
\label{p:eq:rea:product-star}
\ell_m(B_1^\star, \ldots, B_m^\star)
  = (-1)^{m + \sum_i \abs{B_i} + \sum_{i < j} \abs{B_i} \abs{B_j}} \, \ell_m(B_1, \ldots, B_m)^\star.
\end{equation}
On even arguments the sign is $(-1)^m$, and with one odd argument $(-1)^{m + 1}$.

\subparagraph{Reality of the string field.}

The reality condition of the string field~\cite[eq.~(3.6)]{Zwiebach:1993:ClosedStringField}, $\adj{\Psi} = - \bra{\Psi}$, reads $\Psi^\star = - \Psi$ by~\eqref{p:eq:rea:star-bras}.
On the tachyon component
\begin{equation}
\int \frac{\dd^d p}{(2 \pi)^d} \, \phi(p) \, c_1 \bar c_1 \ket{0, p},
\end{equation}
it is the reality condition of a scalar field, $\phi(- p)^* = \phi(p)$, the minus sign coming from the reversal of the two odd modes $c_1 \bar c_1$.
The condition $\Lambda^\star = - \Lambda$ is the one under which every gauge transformation~\eqref{p:eq:conv:gauge} preserves $\Psi^\star = - \Psi$: the free term by the star of $Q_B$,~\eqref{p:eq:conv:QB-star} with $\abs{\Lambda} = 1$, which gives $(Q_B \Lambda)^\star = Q_B \Lambda^\star$ and fixes $\Lambda^\star = - \Lambda$ up to the kernel of $Q_B$, and the interaction terms by the reality of the vertices~\eqref{p:eq:conv:vertex-reality}, through the product rule~\eqref{p:eq:rea:product-star} with one odd argument.

\subparagraph{Reality of the auxiliary fields.}

Every vertex term of the cubic action is a real multiple of a vertex on a tensor whose $N$ strands all have ghost number two.
On even arguments, the reality of the vertices carries the sign $(-1)^N$.
On the slice $\Psi^\star = - \Psi$, $\Sigma_k^\star = (-1)^k \, \Sigma_k$, the strand signs multiply to $(-1)^N$ as well, and we get:
\begin{equation}
\label{p:eq:rea:cubic-real}
\begin{aligned}
  \vertex{N}(\Sigma_k, \Sigma_\ell, \Psi^n)^*
    &
    = (-1)^N \, \vertex{N}\big( \Sigma_k^\star, \Sigma_\ell^\star, (\Psi^\star)^n \big)
  \\
    &
    = (-1)^{N + k + \ell + n} \, \vertex{N}(\Sigma_k, \Sigma_\ell, \Psi^n)
    = \vertex{N}(\Sigma_k, \Sigma_\ell, \Psi^n),
\end{aligned}
\end{equation}
with $N = k + \ell + n$.
The nonpolynomial action is the case without auxiliary fields, $N = n$.
The cubic action is thus real on the slice~\eqref{p:eq:conv:reality-fields} because every strand of ghost number two is star-odd: the field replacing $\Psi^k$ must carry the sign $(-1)^k$ of $(\Psi^\star)^k$.
The sign in the reality condition of $\Sigma_k$ therefore alternates with $k$, and $\Sigma_k$ is not real in a uniform sense.

The source of $\Sigma_k$ is linear in $\Sigma_k$, with a real coefficient once the lower fields are fixed.
It is therefore real exactly when $\Sigma_k$ carries the sign $(-1)^k$ on every direction the source reaches.
For generic vertices, the sources reach every direction and the sign is fixed on the whole field.
For the vertices of a given conformal field theory, the source map need not be surjective.
The directions of $\Sigma_k$ it misses are absent from the dynamics, and on them the condition~\eqref{p:eq:conv:reality-fields} is a choice.

The conditions~\eqref{p:eq:conv:reality-fields} collect these results.
They follow one rule when read on strands: the string field $\Psi$ and the gauge parameter $\Lambda$ are star-odd, and a tensor is star-odd exactly when it has an odd number of star-odd strands.
The strands of ghost number three that occur, those of the outputs of the products on the slice, are star-even, by the same stripping argument as for $\ell_m$.

\subparagraph{Reality of the solution.}

We eliminate the auxiliary fields by inverting their dressed kinetic operators on the sources.
A map built from the vertices commutes with the star on the real slice, by the reality of the vertices, and so does its inverse on a star-stable quotient.
The solution therefore satisfies~\eqref{p:eq:conv:reality-fields} level by level, as a formal series in $\gs$ and wherever the dressed operators are invertible.
Where one of them is degenerate and the source lies in its range, the class of the solution in the quotient is real.
\Cref{p:sec:hs:quintic} works this out at the first order where an inverse appears, with the kernel of $\sprod{2}{2}$ quotiented.


\paragraph*{Reading note.}
The next four sections come from an older companion note on the pairing and the grading at the quartic order.
They derive the same two-string exchange sign by a direct computation on product states (\cref{pg:sec:pairing}), explain why one degree shift, not two, turns the two-string pairing into a graded-antisymmetric form (\cref{pg:sec:grading}), derive the reality of the $\Sigma$ sector of the quartic action from the reality of the quartic vertex (\cref{pg:sec:reality}), and relate the suspended gradings of the homotopy-algebra literature to the physical one (\cref{pg:sec:suspension}).
\status{Verified} The exchange sign and the count of shifts are the result on the two-string pairing quoted above.\record{002-omega2-construction (R7)}
The note uses an earlier notation and a different reality conjugation, listed in the box below.

\begin{quote}
  \small
  \textbf{Translation box: the companion note and the paper.}

  \begin{itemize}
    \setlength{\itemsep}{0.5ex}

    \item $\mc H_0$, the level-matched space at every ghost number, is the space written $\Hs$ in the paper.
      $\epsilon(A) = N_{\rm gh}(A) \bmod 2$ is the parity written $\abs{A}$.

    \item $\braket{A, B}$ and $\braket{X, Y}_2$ are the pairings $\pair{A}{B}$ and $\pair[2]{X}{Y}$, with the same definitions and the same Koszul sign.

    \item $\Sym^2 \mc H_0$ is $\Hn{2}$; the field $\Sigma \in \Sym^2 \mc H_0 / \ker \mc V_4$ is $\Sigma_2 \in \Hn{2} / \ker \sprod{2}{2}$.
      The symbols $\mc V_n$ and $p_{2,2}$ are those of the paper.

    \item The note imposes reality with a conjugation $\ddagger$ for which every field is real: $\Psi^\ddagger = \Psi$, $\Sigma^\ddagger = \Sigma$, and every vertex is equivariant without a sign, $\mc V_n(A_1, \ldots, A_n)^* = \mc V_n(A_1^\ddagger, \ldots, A_n^\ddagger)$.
      On a one-string state of ghost number $N_{\rm gh}$ it is related to the star conjugation of the paper by
      \begin{equation*}
        A^\ddagger
          = (-1)^{N_{\rm gh}(A) + 1} \, A^\star,
      \end{equation*}
      and both act strand by strand on tensors.
      On a $k$-string tensor whose strands have ghost number two, $\ddagger = (-1)^k \star$.
      Hence $\Psi^\ddagger = \Psi$ is $\Psi^\star = - \Psi$, and $\Sigma_k^\ddagger = \Sigma_k$ is $\Sigma_k^\star = (-1)^k \, \Sigma_k$: the two conventions agree for even $k$, and \emph{for odd $k$ the two conditions differ by a sign}.
      A reader who reads $\ddagger$ as $\star$ gets the sign of every odd-$k$ condition wrong.
      The verbal description of $\ddagger$ in \cref{pg:sec:reality:conj}, as the composition of BPZ and Hermitian conjugation, leaves this sign implicit.

    \item The note took the exchange sign of the one-string pairing to be graded antisymmetric.
      The sign derived in the paper's conventions makes it symmetric on its support, and the block below carries the corrected sign as a parameter $s = +1$.
      The two-string results depend on $s^2$ only and are unchanged.
      The graded-antisymmetric form obtained by one shift is even, not odd, and the block says so.

    \item The note writes $g_s$ for $\gs$ and $N_{\rm gh}$ for $\Ngh$; its $\ell_n$ and ghost-number assignments are those of the paper.
  \end{itemize}
\end{quote}

\subsection{The two-string pairing}
\label{pg:sec:pairing}

The kinetic term of $\Sigma$ is a bilinear form on two-string states.
This section builds the pairing that carries it out of the one-string pairing, and computes its exchange sign.
The result is that the two-string pairing is graded \emph{symmetric} in the physical grading, and that the exchange sign of the one-string pairing enters it squared, so that the result does not depend on that sign.
We first record the exchange sign of the one-string pairing, since the two-string computation reuses it slot by slot.

\subsubsection{The one-string pairing and its exchange sign}
\label{pg:sec:pairing:one}

The exchange sign of the one-string pairing is derived in the conventions of the paper, \eqref{p:eq:conv:pairing-symmetry}.
On the support of the pairing we write it as
\begin{equation}
  \label{pg:eq:pairing:one-antisym}
  \braket{A, B}
    = s \, (-1)^{\epsilon(A) \epsilon(B)} \, \braket{B, A},
  \qquad
  s
    = + 1,
\end{equation}
and we keep the sign $s$ explicit below, to show that the two-string result does not depend on it.
Since $\epsilon(A) \epsilon(B) = 0$ on the support, the one-string pairing is symmetric there.
It pairs states of opposite parity: on its support $N_{\rm gh}(A)+N_{\rm gh}(B)=5$, hence
\begin{equation}
  \label{pg:eq:pairing:one-support}
  \epsilon(A) + \epsilon(B)
    \equiv 1 \pmod 2 .
\end{equation}
In particular $\braket{\Psi,\Psi}=0$: the string field at $N_{\rm gh}=2$ pairs with the $N_{\rm gh}=3$ sector.

\subsubsection{Two copies}
\label{pg:sec:pairing:two}

For product states $X=x_1\otimes x_2$ and $Y=y_1\otimes y_2$ in $\mc H_0^{\otimes2}$ we define the two-string pairing as the tensor square of the one-string pairing,
\begin{equation}
  \label{pg:eq:pairing:two-def}
  \boxed{\;
  \braket{X, Y}_2
    := (-1)^{\epsilon(x_2) \epsilon(y_1)} \,
      \braket{x_1 | c_0^- | y_1} \, \braket{x_2 | c_0^- | y_2}
  \;}
\end{equation}
extended bilinearly.
The sign is the Koszul sign for moving $x_2$ past $y_1$, so that the first copy of the pairing acts on the first slots and the second copy on the second slots; with the two $c_0^-$ written out, the pairing is $\braket{\cdot|c_0^-|\cdot}\otimes\braket{\cdot|c_0^-|\cdot}$.
It is non-zero only when both slots saturate, so on its support
\begin{equation}
  \label{pg:eq:pairing:two-support}
  \epsilon(y_i)
    = \epsilon(x_i) + 1
  \quad (i = 1, 2),
  \qquad\text{hence}\qquad
  \epsilon(Y)
    = \epsilon(X),
\end{equation}
where $\epsilon(X)=\epsilon(x_1)+\epsilon(x_2)$: opposite parity in each slot, same total parity.
Off the support both $\braket{X,Y}_2$ and $\braket{Y,X}_2$ vanish, so the exchange relation derived below holds trivially there and we only need to consider the support.

\subsubsection{The exchange sign}
\label{pg:sec:pairing:sign}

Let's exchange $X$ and $Y$ in \eqref{pg:eq:pairing:two-def} and apply \eqref{pg:eq:pairing:one-antisym} in each slot:
\begin{align}
  \braket{Y, X}_2
    &
    = (-1)^{\epsilon(y_2) \epsilon(x_1)} \,
      \braket{y_1, x_1} \, \braket{y_2, x_2}
  \notag
  \\ &
    = (-1)^{\epsilon(y_2) \epsilon(x_1)} \,
      \Big[ s \, (-1)^{\epsilon(x_1) \epsilon(y_1)} \braket{x_1, y_1} \Big]
      \Big[ s \, (-1)^{\epsilon(x_2) \epsilon(y_2)} \braket{x_2, y_2} \Big]
  \notag
  \\ &
    = s^2 \, (-1)^{\epsilon(y_2) \epsilon(x_1) + \epsilon(x_1) \epsilon(y_1) + \epsilon(x_2) \epsilon(y_2)} \,
      \braket{x_1, y_1} \, \braket{x_2, y_2}.
  \label{pg:eq:pairing:exchange-raw}
\end{align}
The two factors $s$ have multiplied to $s^2 = 1$, whichever sign the one-string pairing carries.
It is not yet the exchange sign of the two-string pairing, which still has to be compared with the Koszul sign of \eqref{pg:eq:pairing:two-def}.
Dividing \eqref{pg:eq:pairing:exchange-raw} by \eqref{pg:eq:pairing:two-def},
\begin{equation}
  \braket{Y, X}_2
    = (-1)^{\Xi} \, \braket{X, Y}_2,
  \qquad
  \Xi
    = \epsilon(y_2) \epsilon(x_1)
      + \epsilon(x_1) \epsilon(y_1)
      + \epsilon(x_2) \epsilon(y_2)
      - \epsilon(x_2) \epsilon(y_1).
\end{equation}
On the support we substitute $\epsilon(y_i)=\epsilon(x_i)+1$ from \eqref{pg:eq:pairing:two-support}.
Writing $a_i=\epsilon(x_i)$ and using $a_i^2=a_i$ modulo two,
\begin{equation}
  \begin{aligned}
    \epsilon(y_2) \epsilon(x_1)
      &
      = (a_2 + 1) a_1
      = a_1 a_2 + a_1,
    \\
    \epsilon(x_1) \epsilon(y_1)
      &
      = a_1 (a_1 + 1)
      = 2 a_1
      \equiv 0,
    \\
    \epsilon(x_2) \epsilon(y_2)
      &
      = a_2 (a_2 + 1)
      \equiv 0,
    \\
    - \epsilon(x_2) \epsilon(y_1)
      &
      \equiv a_2 (a_1 + 1)
      = a_1 a_2 + a_2 .
  \end{aligned}
\end{equation}
The two cross terms $a_1a_2$ cancel, and
\begin{equation}
  \Xi
    \equiv a_1 + a_2
    = \epsilon(X)
    \pmod 2 .
\end{equation}
Since $\epsilon(Y)=\epsilon(X)$ on the support, $\epsilon(X)=\epsilon(X)^2=\epsilon(X)\epsilon(Y)$, and we obtain the exchange relation:
\begin{equation}
  \label{pg:eq:pairing:symmetric}
  \boxed{\;
  \braket{Y, X}_2
    = (-1)^{\epsilon(X) \epsilon(Y)} \, \braket{X, Y}_2 .
  \;}
\end{equation}
The two-string pairing is graded symmetric in the physical grading: $+1$ when at least one argument is even, a genuine $-1$ when both are odd.
The exchange sign of the one-string pairing enters squared, so the result holds for either sign $s$.
The opposite-parity support \eqref{pg:eq:pairing:two-support} only serves to cancel the cross terms; the same symmetric verdict is obtained on a model with even saturation, so it does not hinge on the precise value of the sphere anomaly.

Two further facts about $\braket{\cdot,\cdot}_2$ are used later.
It is non-degenerate whenever the one-string pairing is, on a finite-dimensional model, where $\det\braket{\cdot,\cdot}_2=(\det\braket{\cdot,\cdot})^{2\dim\mc H_0}$; on the level-matched space one expects the same to follow from the non-degeneracy of $\braket{\cdot,\cdot}$ on $\mc H_0$, but the infinite-dimensional statement is not established.
And it is the pairing fixed by the kinetic term: since $\mc V_4$ and $p_{2,2}$ are related through \eqref{p:eq:conv:sprod}, the kinetic term $\mc V_4(\Sigma,\Sigma)=\braket{\Sigma,p_{2,2}\Sigma}_2$ singles out this pairing as the natural one.
Whether it is the unique pairing with respect to which $p_{2,2}$ is cyclic is not established.

The exchange sign \eqref{pg:eq:pairing:symmetric} and the support statements were checked symbolically on a four-state model with ghost numbers $(2,3,2,3)$, graded by the physical parity only.



\paragraph*{Reading note.}
The next section works out the grading of the two-string pairing.
It shows that the two shifts one would expect from the two tensor factors are invisible to the statistics, and that a single shift of the whole field gives the graded-antisymmetric presentation.

\subsection{One shift, not two}
\label{pg:sec:grading}

The natural expectation runs as follows.
The symplectic form of closed string field theory is graded antisymmetric and odd, and in the homotopy-algebra literature this is arranged by a degree shift of the string field.
The auxiliary field $\Sigma$ has two tensor factors; each factor is a string field; so one expects two degree-shift units, one per factor, in the symplectic partner of $\braket{\cdot,\cdot}_2$.
The correct count is one.
This section writes the two-unit expectation down, shows where it fails, and identifies the single shift that does the job.
The short version is that the \enquote{two} of the two tensor factors is real but counts $c_0^-$ insertions, and that count fixes the operation parity, not the suspension.

\subsubsection{The naive count, and where it fails}
\label{pg:sec:grading:naive}

Let's write the expectation as a formula.
A degree shift $s\colon\mc H_0\to\mc H_0[1]$ applied to each tensor factor changes the parity of a product state by
\begin{equation}
  \epsilon(x_1 \otimes x_2)
    \longmapsto
    \big( \epsilon(x_1) + 1 \big) + \big( \epsilon(x_2) + 1 \big)
    = \epsilon(X) + 2
    \equiv \epsilon(X) \pmod 2 .
\end{equation}
Two shift units add up to an even integer and do nothing to the $\Z_2$ parity of $X$.
A parity relabelling that changes nothing cannot change an exchange sign: after the shift the pairing satisfies exactly the same relation \eqref{pg:eq:pairing:symmetric} as before and is still graded symmetric.
The two-unit count therefore fails, not because the count of factors is wrong, but because an even shift is invisible to the statistics.
The same remark applies to the shift $A_c=s^{-2}\tilde A_c$ by which Münster and Sachs place the closed string field in degree zero \cite{Muenster:2012:HomotopyAlgebrasString-kn,Muenster:2013:HomotopyAlgebrasQuantum}: it is an integer shift by $-2$, even, and it does not touch the parity \eqref{p:eq:conv:parity}.

\subsubsection{The single shift}
\label{pg:sec:grading:single}

Now let's ask directly which relabelling $\epsilon\mapsto\epsilon+c_s$ turns \eqref{pg:eq:pairing:symmetric} into a graded-antisymmetric relation,
\begin{equation}
  \label{pg:eq:grading:antisym-demand}
  \braket{Y, X}_2
    \overset{?}{=}
    - (-1)^{(\epsilon(X) + c_s)(\epsilon(Y) + c_s)} \, \braket{X, Y}_2 .
\end{equation}
Comparing with \eqref{pg:eq:pairing:symmetric}, the requirement is
\begin{equation}
  (-1)^{\epsilon(X) \epsilon(Y)}
    = - (-1)^{(\epsilon(X) + c_s)(\epsilon(Y) + c_s)}
  \quad\text{on the support.}
\end{equation}
With $a=\epsilon(X)$, $b=\epsilon(Y)$ and $a=b$ on the support \eqref{pg:eq:pairing:two-support}, the exponent on the right is $(a+c_s)^2\equiv a+c_s$, and the left exponent is $a^2\equiv a$; the relation holds if and only if $c_s\equiv1$.
So exactly one odd shift unit is required, and one is consistent:
\begin{equation}
  \label{pg:eq:grading:symp}
  \boxed{\;
  \begin{aligned}
    \braket{\cdot, \cdot}_2^{\rm symp}
      &
      := \braket{\cdot, \cdot}_2 \circ (s \otimes s)
      \quad\text{on}\quad
      \big( \mc H_0^{\otimes 2} \big)[1],
    \\
    \braket{Y, X}_2^{\rm symp}
      &
      = - (-1)^{\epsilon_s(X) \epsilon_s(Y)} \, \braket{X, Y}_2^{\rm symp},
  \end{aligned}
  \;}
\end{equation}
where $\epsilon_s=\epsilon+1$ is the shifted parity.
The notation $s\otimes s$ deserves one sentence, because it is the trap: the two factors of $s\otimes s$ act on the two \emph{arguments} $X$ and $Y$ of the bilinear form, each argument being one element of $\mc H_0^{\otimes2}$ shifted once as a whole.
They do not act on the two tensor factors $x_1$, $x_2$ of a single argument.
The shift is per field, not per strand, because $\Sigma$ is one element of the enlarged complex $\Sym^2\mc H_0$, one entry of the form, hence one $C\to C[1]$, irrespective of how many single-string factors build it.

\subsubsection{The two counts reconciled}
\label{pg:sec:grading:reconcile}

The two $\Z_2$ data are now separated.
The number of $c_0^-$ insertions, two, is the number of tensor factors of $\Sigma$; it fixes the operation parity of $\braket{\cdot,\cdot}_2$, even, hence its graded symmetry.
The number of suspension units, one, is the number of fields; it is the shift that produces the graded-antisymmetric presentation \eqref{pg:eq:grading:symp}.
This form is even, not odd: $\braket{\cdot,\cdot}_2$ pairs states of equal total parity, so it is even in every grading whose parity is a function of the ghost number, and an odd pairing for $\Sigma$ needs a grading of another kind.
The expected \enquote{two} is real, but it makes the pairing symmetric, the opposite of what a naive count suggests.
The shift itself is a choice of presentation, not a fact about $\Sigma$: in the physical grading the two-string pairing is graded symmetric, full stop, and the antisymmetric object exists only in the shifted presentation.
Nothing downstream depends on the label: the entries of $\braket{\cdot,\cdot}_2$, of $\mc V_4$ and of $p_{2,2}$ are the same in both presentations, and so are the kernel of $p_{2,2}$ and its cyclicity.

\subsubsection{The smallest instance}
\label{pg:sec:grading:example}

Take four states of definite ghost number, $x_1,x_2$ at $N_{\rm gh}=2$ and $y_1,y_2$ at $N_{\rm gh}=3$, so that each slot saturates.
Then $\epsilon(x_i)=0$, $\epsilon(y_i)=1$, $\epsilon(X)=\epsilon(Y)=0$.
The Koszul sign in \eqref{pg:eq:pairing:two-def} is $(-1)^{\epsilon(x_2)\epsilon(y_1)}=(-1)^{0}=+1$, and
\begin{align}
  \braket{X, Y}_2
    &
    = \braket{x_1, y_1} \braket{x_2, y_2},
  \\
  \braket{Y, X}_2
    &
    = (-1)^{\epsilon(y_2)\epsilon(x_1)} \braket{y_1, x_1} \braket{y_2, x_2}
    = s^2 \braket{x_1, y_1} \braket{x_2, y_2}
    = + \braket{X, Y}_2,
\end{align}
using $\braket{y_i,x_i}=s\,(-1)^{0}\braket{x_i,y_i}=s\braket{x_i,y_i}$ in each slot, with $s = +1$ from \eqref{pg:eq:pairing:one-antisym}.
The prediction of \eqref{pg:eq:pairing:symmetric} is $(-1)^{0\cdot0}=+1$: correct.
Graded antisymmetry in the physical grading would predict $-(-1)^{0}=-1$: wrong.
With one shift, $\epsilon_s(X)=\epsilon_s(Y)=1$, the shifted rule \eqref{pg:eq:grading:symp} predicts $-(-1)^{1\cdot1}=+1$: correct again.

For the other sign, take $x_1$ at $N_{\rm gh}=2$, $x_2$ at $3$, $y_1$ at $3$, $y_2$ at $2$, so that $\epsilon(X)=\epsilon(Y)=1$.
Now the Koszul sign is $(-1)^{\epsilon(x_2)\epsilon(y_1)}=(-1)^{1}=-1$, and
\begin{align}
  \braket{X, Y}_2
    &
    = - \braket{x_1, y_1} \braket{x_2, y_2},
  \\
  \braket{Y, X}_2
    &
    = (-1)^{\epsilon(y_2)\epsilon(x_1)} \braket{y_1, x_1} \braket{y_2, x_2}
    = (+1) \, s^2 \braket{x_1, y_1} \braket{x_2, y_2}
    = - \braket{X, Y}_2,
\end{align}
since $\epsilon(x_1)\epsilon(y_1)=0\cdot1=0$ and $\epsilon(x_2)\epsilon(y_2)=1\cdot0=0$ in the two slot exchanges.
The prediction of \eqref{pg:eq:pairing:symmetric} is $(-1)^{1\cdot1}=-1$: correct, and a genuine minus sign on the odd--odd sector.
Graded antisymmetry would predict $-(-1)^{1}=+1$: wrong.
With one shift, $\epsilon_s(X)=\epsilon_s(Y)=0$, the shifted rule predicts $-(-1)^{0}=-1$: correct.
The two instances together show that $c_s=1$ is not fitted to one sector: it flips the relation consistently on both.



\paragraph*{Reading note.}
\status{mixed} The next section is the reality argument of the companion note at the quartic order, in the $\ddagger$ convention of the translation box above.\record{028-reality-tower (R35)}
It is the case $k = 2$ of the reality of the auxiliary tower, conditional on the equivariance of the quartic vertex, which the section states as a hypothesis where it is used.

\subsection{Reality}
\label{pg:sec:reality}

The action must be real.
For the $\Psi$ sector this is standard: with the closed-string reality condition $\Psi^\ddagger=\Psi$, the kinetic term $\frac12\braket{\Psi|c_0^-Q_B|\Psi}$ is real \cite{Erbin:2021:StringFieldTheory}.
This section does the same for the $\Sigma$ sector of the quartic action~\eqref{p:eq:hs4:S4tilde}: the conjugation of the pairing, of the vertex and of $p_{2,2}$, the condition on $\Sigma$ that follows, and the sense in which that condition is forced rather than chosen.
The argument is algebraic and rests on one hypothesis about the vertex, stated where it is used.

\subsubsection{The conjugation}
\label{pg:sec:reality:conj}

The conjugation $\ddagger$ is the composition of BPZ and Hermitian conjugation.
It is an antilinear involution of $\mc H_0$ preserving the ghost number: $(aA+bB)^\ddagger=a^*A^\ddagger+b^*B^\ddagger$ and $(A^\ddagger)^\ddagger=A$.
A state is \emph{real} if $A^\ddagger=A$; every state splits as $A=A_r+\I A_i$ with $A_r=\frac12(A+A^\ddagger)$ and $A_i=\frac1{2\I}(A-A^\ddagger)$ both real, so the real states span $\mc H_0$ over $\C$.
On two-string states we extend it componentwise,
\begin{equation}
  (A_1 \otimes A_2)^\ddagger
    := A_1^\ddagger \otimes A_2^\ddagger,
\end{equation}
with no reordering sign.
For $\Sigma$ this is unambiguous: both strands sit at $N_{\rm gh}=2$, even, so the alternative convention with reversed order, which differs by $(-1)^{k(k-1)/2}$ on tensors all of whose strands are odd, agrees with it.
$\Sym^2\mc H_0$ is $\ddagger$-stable, and $(\Psi^2)^\ddagger=\Psi^\ddagger\otimes\Psi^\ddagger=\Psi^2$.

A multilinear form $F$ is \emph{real} if it takes real values on real arguments.
This is equivalent to \emph{equivariance},
\begin{equation}
  \label{pg:eq:reality:equiv}
  F(A_1, \ldots, A_n)^*
    = F(A_1^\ddagger, \ldots, A_n^\ddagger)
  \quad\text{for all } A_i ,
\end{equation}
since both sides of \eqref{pg:eq:reality:equiv} are antilinear in each slot and agree on real arguments, which span.
We use the two formulations interchangeably.

\subsubsection{Conjugation of the ingredients}
\label{pg:sec:reality:ingredients}

\emph{The pairing.}
In the normalization $\braket{c_{-1}\bar c_{-1}c_0^-c_0^+c_1\bar c_1}=-\frac12$, real and without a factor of $\I$, the pairing $\braket{A,B}=\braket{A|c_0^-|B}$ is real on real states.
The only structural input is that Hermitian conjugation of $c_0^-$ and of the BPZ product introduces no phase, which is what the normalization fixes; it is the same choice that makes the $\Psi$ kinetic term real, and the reality of the $\Psi$ kinetic term already pairs $N_{\rm gh}=2$ with $N_{\rm gh}=3$.
One caveat: Hermitian conjugation reverses the order of $A$ and $B$, so for two odd arguments the equivariance of $\braket{\cdot,\cdot}$ may carry a relative sign $(-1)^{\epsilon(A)\epsilon(B)}$ that the even--odd case does not probe.
Every pairing below has one even argument, so the odd--odd case is never needed.
The two-string pairing is then real too: on product states $\braket{X,Y}_2$ is a product of two real pairings, both sides of the equivariance identity are antilinear in each argument, and product states span.

\emph{The vertex.}
The hypothesis the whole section rests on is that the quartic vertex is equivariant,
\begin{equation}
  \label{pg:eq:reality:V4}
  \mc V_4(A_1, A_2, A_3, A_4)^*
    = \mc V_4(A_1^\ddagger, A_2^\ddagger, A_3^\ddagger, A_4^\ddagger).
\end{equation}
This is the standard reality of a closed-string vertex, the one that makes the non-polynomial action real: it follows from the Hermiticity of the vertex operators and plumbing data, $\mc O(t^*)=\mc O(t)^*$ in the language of Sonoda and Zwiebach \cite{Sonoda:1990:CovariantClosedString}, together with the invariance of the moduli region of the vertex and of its measure under complex conjugation of the moduli.
We take \eqref{pg:eq:reality:V4} from the literature and do not re-derive it; everything below is conditional on it.

\emph{The product.}
Since $\braket{\cdot,\cdot}_2$ is real and non-degenerate, \eqref{pg:eq:reality:V4} says that $p_{2,2}$ is a real operator, $p_{2,2}(Y^\ddagger)=(p_{2,2}Y)^\ddagger$.
Indeed, conjugating \eqref{p:eq:conv:sprod} in two ways,
\begin{equation}
  \begin{aligned}
    \braket{X^\ddagger, (p_{2,2} Y)^\ddagger}_2
      &
      = \braket{X, p_{2,2} Y}_2^*
      = \mc V_4(X, Y)^*
    \\ &
      = \mc V_4(X^\ddagger, Y^\ddagger)
      = \braket{X^\ddagger, p_{2,2}(Y^\ddagger)}_2
    \quad\text{for all } X,
  \end{aligned}
\end{equation}
and non-degeneracy gives the claim.
A real operator has a $\ddagger$-stable kernel and image, and its inverse, where it exists, is real; in particular $\ker\mc V_4=\ker p_{2,2}$ is $\ddagger$-stable, so the quotient~\eqref{p:eq:hs4:quotient} inherits the conjugation and a reality condition on the class of $\Sigma$ makes sense.

\subsubsection{The reality condition on \texorpdfstring{$\Sigma$}{Sigma}}
\label{pg:sec:reality:condition}

Impose
\begin{equation}
  \label{pg:eq:reality:Sigma}
  \boxed{\;
  \Sigma^\ddagger = \Sigma .
  \;}
\end{equation}
Then both $\Sigma$ terms of \eqref{p:eq:hs4:S4tilde} are real: by \eqref{pg:eq:reality:V4}, $\mc V_4(\Sigma,\Sigma)^*=\mc V_4(\Sigma^\ddagger,\Sigma^\ddagger)=\mc V_4(\Sigma,\Sigma)$, and $\mc V_4(\Sigma,\Psi^2)^*=\mc V_4(\Sigma,\Psi^2)$ likewise, the coefficients being real.
The on-shell value $\Sigma=\Psi^2$ is real, and the class of a real solution in the quotient is real.
Reality is thus conditional on \eqref{pg:eq:reality:V4} and on nothing else; the two $c_0^-$ of $\braket{\cdot,\cdot}_2$ play no role beyond the reality of the pairing, and if a convention had attached a phase $\lambda$ to the single $c_0^-$ of $\braket{\cdot,\cdot}$, it would enter the $\Sigma$ kinetic term as $\lambda^2$.

Let's now ask whether \eqref{pg:eq:reality:Sigma} is forced.
The kinetic term alone does not force it: for $\Sigma=\I\Sigma_r$ with $\Sigma_r$ real, $\mc V_4(\Sigma,\Sigma)=-\mc V_4(\Sigma_r,\Sigma_r)$ is real as well, so $\Sigma^\ddagger=-\Sigma$ would also do.
The source term removes this freedom.
It is linear in $\Sigma$, and for $\Sigma=\Sigma_r+\I\Sigma_i$,
\begin{equation}
  \mc V_4(\Sigma, \Psi^2)
    = \mc V_4(\Sigma_r, \Psi^2)
      + \I \, \mc V_4(\Sigma_i, \Psi^2),
\end{equation}
which is real for all real $\Psi$ if and only if $\mc V_4(\Sigma_i,\Psi^2)=0$ for all real $\Psi$.
For a generic symmetric vertex the covectors $\mc V_4(\,\cdot\,,\Psi^2)$ span the dual of $\Sym^2\mc H_0$ and $\Sigma_i=0$; for the actual vertex, the directions annihilated by every $\mc V_4(\,\cdot\,,\Psi^2)$ are directions absent from the dynamics, and \eqref{pg:eq:reality:Sigma} is forced up to them.
A relative phase $\Sigma^\ddagger=\e^{\I\theta}\Sigma$ is excluded in the same way, and $\Sigma\to-\Sigma$ is a field redefinition that flips the sign of the source term, not a different reality condition.
So \eqref{pg:eq:reality:Sigma} is the unique condition compatible with a real $\Psi$, which is how the quartic step reads it off the equation of motion.

The argument was checked on a finite-dimensional model with a non-trivial antilinear involution, a random real reflection composed with complex conjugation, so that the real subspace is not the coordinate one, including a degenerate kernel and the two failure modes (a non-real component in a vertex, a relative phase on $\Sigma$).



\paragraph*{Reading note.}
The last section of the companion note relates the suspended gradings of Cieliebak, Fukaya and Latschev and of Münster and Sachs to the physical one.

\subsection{Suspended versus physical grading}
\label{pg:sec:suspension}

The two frameworks from which the algebraic vocabulary of this construction is borrowed both suspend, and both fold the load-bearing sign of the odd symplectic form into the shift.
This section says what each one does and how to undo it for the objects of this note.

\emph{Cieliebak, Fukaya and Latschev} \cite{Cieliebak:2015:HomologicalAlgebraRelated} build their symmetric algebras on the suspension $C[1]$ of the underlying complex, so their degree is the physical degree minus one, $|c|=\deg c-1$.
Their structure maps carry a degree $-2d(k+g-1)-1$ depending on the numbers $k$ of inputs and $g$ of handles and on a fixed integer $d$, but not on the number of outputs; the grading is homological, the differential lowers the degree, opposite to $Q_B$ which raises the ghost number.

\emph{Münster and Sachs} \cite{Muenster:2012:HomotopyAlgebrasString-kn,Muenster:2013:HomotopyAlgebrasQuantum} place the closed string field in degree zero after the shift $A_c=s^{-2}\tilde A_c$, and write the symplectic form as
\begin{equation}
  \omega_c
    = (\,\cdot\,, c_0^- \,\cdot\,)_{\rm bpz} \circ (s^2 \otimes s^2),
\end{equation}
which is odd and graded antisymmetric by construction of the shift: the apparent symmetry of the BPZ product is removed by the suspension.
The physical ghost number of $\tilde A_c$ is kept implicit.

\emph{Undoing it.}
For the objects at this order the translation is the following.
The physical grading is $\epsilon=N_{\rm gh}\bmod2$ with $\Psi$ at $2$ and $\Sigma$ at $4$, both even.
The one-string pairing is symmetric on its support (\cref{pg:sec:pairing:one}), with no shift needed; the shift $s^{-2}$ is even and does not touch the parity.
The two-string pairing is graded symmetric (\cref{pg:sec:pairing:sign}); its graded-antisymmetric partner, an even form, is the one-unit shift \eqref{pg:eq:grading:symp}, and that shift is a choice of presentation, not a property of $\Sigma$.
The ghost numbers of the maps are fixed by saturation: $N_{\rm gh}(\ell_n)=3-2n$ and $N_{\rm gh}(p_{2,2})=+2$ from \eqref{p:eq:conv:ghost-numbers}, every output at $N_{\rm gh}=3$ and every input at $N_{\rm gh}=2$.
The entries of $\braket{\cdot,\cdot}_2$, $\mc V_4$ and $p_{2,2}$ are the same in the suspended and in the physical presentation; only Koszul-sign labels differ, and in particular the subspace $\wedge^2\mc H_0$ of antisymmetric two-string states, which lies in the kernel~\eqref{p:eq:hs4:kernel} of $p_{2,2}$ by the symmetry of the vertex and decouples from the action, is the same subspace in both, even though after a one-unit shift per strand it would be called the graded-symmetric square.
A reader working in the suspended presentation must select the subspace, not the label.
The same translation, carried through every relation of the bracket-and-cobracket structure at higher order, is done in the chapter on the IBL structure (\cref{ch:ibl}).




\subsection{The two normalizations of the sphere correlator}
\label{sec:conv:two-normalizations}

\status{Verified}\record{041-convention-reconciliation (R43)}\record{046-stage0-translation-check}\footnote{Derivation record not published (it holds data of the companion $\group{SL}(2, \C)$ project).}

Two normalizations of the ghost correlator on the sphere are met in the literature.
The paper and this write-up follow~\cite{Erbin:2021:StringFieldTheory}, with~\eqref{p:eq:conv:sphere}:
\begin{equation}
  \label{eq:conv:nu-book}
  \nu
    := \braket{c_{-1} \bar c_{-1} c_0^- c_0^+ c_1 \bar c_1}
    = - \frac{1}{2}.
\end{equation}
The working normalization of Belopolsky and Zwiebach~\cite{Belopolsky:1995:OffshellClosedString} is $\nu = -1$, and it is the one in which the explicit quartic vertex with $\group{SL}(2, \C)$ local coordinates of \cref{ch:sl2c-vertex} was evaluated.
This section shows that the two differ only by the normalization of the sphere vacuum, gives the conversion of every quantity used in the write-up, and states what the checks establish.

\paragraph{The zero modes are the same.}
The two normalizations use the same zero modes $c_0^\pm = \tfrac{1}{2} (c_0 \pm \bar c_0)$ and $b_0^\pm = b_0 \pm \bar b_0$, and the same ordering $c_0^- c_0^+$.
Since $c_0^-$ and $c_0^+$ anticommute, the product of the two is
\begin{equation}
  c_0^- c_0^+
    = \frac{1}{4} \, (c_0 - \bar c_0) (c_0 + \bar c_0)
    = \frac{1}{4} \, (c_0 \bar c_0 - \bar c_0 c_0)
    = \frac{1}{2} \, c_0 \bar c_0,
\end{equation}
in both, so that $\nu = \tfrac{1}{2} \braket{c_{-1} \bar c_{-1} c_0 \bar c_0 c_1 \bar c_1}$ in both.
The opposite ordering $c_0^+ c_0^-$ would flip the sign of $\nu$, not its magnitude.
The difference between $\nu = -1$ and $\nu = -\tfrac{1}{2}$ is therefore not a convention on the zero modes: it is the value of the base correlator, $\braket{c_{-1} \bar c_{-1} c_0 \bar c_0 c_1 \bar c_1} = -2$ against $-1$, that is the normalization of the sphere vacuum, with the same sign.

\paragraph{The conversion.}
Every sphere correlator computed with $\nu = -1$ is twice the same correlator computed with~\eqref{eq:conv:nu-book}.
A quantity built from a net number $k$ of sphere correlators (negative when it contains an inverse) is therefore converted to the normalization of this write-up by the factor $2^{-k}$, and its sign is unchanged:
\begin{itemize}
  \setlength{\itemsep}{0.5ex}

  \item the pairing, one correlator: $\times \tfrac{1}{2}$; for instance $\braket{T | c_0^- c_0^+ | T} = \nu$ for the tachyon $T = c_1 \bar c_1 \ket{0}$, $-\tfrac{1}{2}$ against $-1$;

  \item the quartic vertex $\vertex{4}$ on two-string states, one four-punctured sphere although its arguments are two-string states: $\times \tfrac{1}{2}$;

  \item the matrix of the two-string pairing, a product of two pairings: $\times \tfrac{1}{4}$;

  \item the operator $\sprod{2}{2}$, which is the quartic vertex read through the inverse of the two-string pairing, and its eigenvalues: $\times 2$;

  \item the constants of the three- and four-point tachyon couplings: $\times \tfrac{1}{2}$ each.
\end{itemize}
The last item is a check of the rule rather than a consequence: the cubic tachyon coupling was recomputed from~\eqref{eq:conv:nu-book} and the mapping radii of the cubic vertex, without using the factor $\tfrac{1}{2}$.
It agrees with the rule: equal signs, magnitude halved.
The constant of the quartic contact term, quoted on the potential side, differs by the same factor two; the value in this normalization is fixed by the cubic vertex, the pairing and factorization; its sign depends on the orientation of the measure, \cref{sec:conv:measure}.

\paragraph{The same theory.}
The conversion is a rescaling of the field and of the coupling.
Each term of the action~\eqref{p:eq:conv:action} is one sphere correlator, so the action computed with $\nu = -1$ is twice the action of this write-up at the same field and coupling.
Writing $S_{\nu}[\Psi; \gs]$ for the action in the normalization $\nu$, the term with $n$ fields scales as $\gs^{n-2} \, \Psi^n$, and the factor two is absorbed by
\begin{equation}
  \label{eq:conv:rescaling}
  S_{-1/2}[\Psi; \gs]
    = S_{-1} \big[ \Psi / \sqrt{2}; \sqrt{2} \, \gs \big],
\end{equation}
since $2 \, (\sqrt{2} \, \gs)^{n-2} \, (\Psi / \sqrt{2})^n = \gs^{n-2} \, \Psi^n$ for every $n$.
The field of this write-up is $\sqrt{2}$ times the field of the other normalization, and the coupling of the other normalization is $\sqrt{2}$ times the coupling here.
Everything a field and coupling rescaling leaves alone is the same in the two normalizations: the sign of $\nu$, the tachyon mass $m_T^2 = - 4 / \alpha'$, the rank of $\vertex{4}$, the signature of the two-string pairing, the ratios of eigenvalues, and the zeros of the determinant of $\vertex{4}$ as a function of the total momentum.
Only magnitudes are converted, and the conversion carries no physics.

\paragraph{What is established.}
The sign of the conversion is established by two independent routes: the normalization of~\cite{Erbin:2021:StringFieldTheory} with the mode algebra, and the kinetic term of the tachyon, which uses no vacuum normalization.
The magnitude $\tfrac{1}{2}$ is established by the recomputation of the cubic coupling above.
That no further factor enters (a different normalization of the matter correlator, a volume factor, a normalization of the coupling) is established up to four points by the uniform ratio two of the pairing, the cubic coupling and the quartic contact constant, which come from separately published constants.
What remains a definition, and no computation can change, is the factor $\tfrac{1}{2}$ as the ratio of the two vacuum normalizations.
The conversion is checked by the companion script on conventions, listed in the script table of \cref{ch:finite-models}.
The orientation of the measure on the moduli space is a separate question, which a three-point function cannot probe; it is the subject of the next section.


\subsection{Orientation of the off-shell measure}
\label{sec:conv:measure}

\status{Proven}\record{042-q2-measure-orientation (R47)}

Two transcriptions of the genus-zero measure prefactor are in use.
Belopolsky and Zwiebach write a factor $\I / 2\pi$ per complex modulus~\cite{Belopolsky:1995:OffshellClosedString}, while~\cite{Erbin:2021:StringFieldTheory} writes $(2\pi \I)^{-1} = - \I / 2\pi$.
Read side by side, they differ by $(-1)^{N-3}$ for an $N$-point amplitude.
A conflict between them would be physical: the ratio $C_4 / C_3^2$ of the quartic and cubic coupling constants and the sign of the eigenvalues of $\sprod{2}{2}$ flip with it, while the cubic constant is untouched.
This section shows that the two transcriptions do not conflict.
The prefactor enters the amplitude only multiplied by a second sign, the convention for the antiholomorphic contour integral, which neither formula displays; once each prefactor is read with the contour convention of its own source, the two agree.

\paragraph{Genus zero and on-shell states.}
We work at genus zero, with on-shell external states $\mathscr V_i = c \bar c V_i$, $V_i$ a matter primary of weight $(1, 1)$, trivial local coordinates (on shell the correlator does not depend on them), and the conventions of this chapter.
The value of one overall sign, stated at the end, is taken as an input from a separate computation.

\paragraph{Orientation conventions.}
These are conventions, free and fixed once.
A complex modulus is $m = m_r + \I m_i$ with $m_r, m_i$ real, and the pair $(m_r, m_i)$ is positively oriented, so that $\dd m \wedge \dd \bar m = - 2 \I \, \dd m_r \wedge \dd m_i$.
The area element is $\dd^2 m := 2 \, \dd m_r \, \dd m_i$, the convention of~\cite{Erbin:2021:StringFieldTheory}.
For several moduli $m^{(1)}, \ldots, m^{(N-3)}$, the positively oriented real frame is $(m^{(1)}_r, m^{(1)}_i, m^{(2)}_r, m^{(2)}_i, \ldots)$, real and imaginary parts kept adjacent.
The contours around the punctures are anticlockwise, with the local coordinate patch on their left, so that $\oint \frac{\dd w}{2\pi \I} \, w^{-1} = 1$.

\paragraph{Two signs.}
We keep two signs free.
The prefactor of the measure is $\kappa / 2\pi$ per complex modulus: Belopolsky and Zwiebach have $\kappa = + \I$, and~\cite{Erbin:2021:StringFieldTheory} has $\kappa = - \I$.
The antiholomorphic contour integral is normalized by a sign $\zeta$,
\begin{equation}
  \label{eq:conv:zeta}
  \oint_C \frac{\dd \bar w}{2\pi \I} \, \bar w^{-n-2}
    =: \zeta \, \delta_{n,-1},
  \qquad
  \zeta
    = \pm 1.
\end{equation}
The prefactor $\kappa$ is a transcription, a claim that can be wrong; $\zeta$ is a convention.
Neither is fitted to a coupling constant: the anchor of the computation is the on-shell amplitude gauge-fixed by the Faddeev--Popov procedure, which is independent of both.

\paragraph{Step 1: the value of $\zeta$ under the anticlockwise rule.}
The book~\cite{Erbin:2021:StringFieldTheory} states that the contours run anticlockwise and writes the antiholomorphic part of the antighost insertion as $\oint_C \frac{\dd \bar z}{2\pi \I} \, \bar b(\bar z) \bar v(\bar z)$ along the same contour $C$.
Taken literally, $\dd \bar w$ is the differential of $\bar w$ along the anticlockwise circle $w = r \e^{\I \theta}$, $\theta \colon 0 \to 2\pi$, so that $\bar w = r \e^{- \I \theta}$ and $\dd \bar w = - \I r \e^{- \I \theta} \dd \theta$.
Then
\begin{equation}
  \label{eq:conv:zeta-literal}
  \oint \frac{\dd \bar w}{2\pi \I} \, \bar w^{-n-2}
    = \frac{1}{2\pi \I} \int_0^{2\pi} (- \I r \e^{- \I \theta}) \, r^{-n-2} \e^{\I (n+2) \theta} \, \dd \theta
    = - \frac{r^{-n-1}}{2\pi} \int_0^{2\pi} \e^{\I (n+1) \theta} \, \dd \theta
    = - \delta_{n,-1},
\end{equation}
and the literal reading gives $\zeta = -1$.
The other common reading treats $\bar w$ as an independent variable with its own anticlockwise contour, so that $\oint \frac{\dd \bar w}{2\pi \I} = + \Res_{\bar w}$, and gives $\zeta = + 1$.
Both are legitimate conventions, and they are not the same one.
The literal reading is the one of~\cite{deLacroix:2017:ClosedSuperstringField}, which states, with the same rule on the orientation of the contours, that integration over antiholomorphic coordinates comes with a factor $(- 2\pi \I)^{-1}$: this is $\zeta = -1$.

\paragraph{Step 2: the antighost insertions of one modulus.}
At four points we fix the punctures $z_1, z_2, z_3$ and take as modulus the position $z_4 = x + \I y$ of the fourth puncture, with local coordinate $F_4(w) = z_4 + w$.
The Schiffer vectors are $\pd F_4 / \pd x = 1$ and $\pd F_4 / \pd y = \I$, with conjugates $1$ and $- \I$.
With $\oint \frac{\dd w}{2\pi \I} \, b(w) = b_{-1}$ and, by~\eqref{eq:conv:zeta}, $\oint \frac{\dd \bar w}{2\pi \I} \, \bar b(\bar w) = \zeta \, \bar b_{-1}$, the two insertions are
\begin{equation}
  \label{eq:conv:Bxy}
  B_x
    = b_{-1} + \zeta \, \bar b_{-1},
  \qquad
  B_y
    = \I \, (b_{-1} - \zeta \, \bar b_{-1}).
\end{equation}

\paragraph{Step 3: their product.}
In the order $B_x B_y$ fixed by the orientation, and with $b_{-1}^2 = \bar b_{-1}^2 = 0$ and $\bar b_{-1} b_{-1} = - b_{-1} \bar b_{-1}$:
\begin{equation}
  \label{eq:conv:BxBy}
  B_x B_y
    = \I \, (b_{-1} + \zeta \, \bar b_{-1}) (b_{-1} - \zeta \, \bar b_{-1})
    = \I \, (- \zeta \, b_{-1} \bar b_{-1} + \zeta \, \bar b_{-1} b_{-1})
    = - 2 \I \zeta \, b_{-1} \bar b_{-1}.
\end{equation}
This $- 2 \I$ is the factor of $\dd m \wedge \dd \bar m = - 2 \I \, \dd m_r \wedge \dd m_i$, appearing on the side of the ghosts instead of the side of the measure.
The order of the two insertions and the sign $\zeta$ enter only through their product, and the sign of the Schiffer vector drops out because it appears twice.

\paragraph{Step 4: the action on the fourth puncture.}
The vertex operator $\mathscr V_4 = c \bar c V_4$ is the state $c_1 \bar c_1 \ket{V_4}$.
With $\{ b_m, c_n \} = \delta_{m+n,0}$, $b_{-1} \ket{V_4} = \bar b_{-1} \ket{V_4} = 0$ and $\bar b_{-1}$ anticommuting with $c_1$,
\begin{equation}
  \label{eq:conv:ghost-strip}
  \bar b_{-1} c_1 \bar c_1 \ket{V_4}
    = - c_1 \bar b_{-1} \bar c_1 \ket{V_4}
    = - c_1 \ket{V_4},
  \qquad
  b_{-1} \bar b_{-1} c_1 \bar c_1 \ket{V_4}
    = - \ket{V_4}.
\end{equation}
The pair $B_x B_y$ and every $\mathscr V_i$ are Grassmann even, so moving the pair onto the fourth puncture costs no sign, and with~\eqref{eq:conv:BxBy}
\begin{equation}
  \label{eq:conv:four-point-correlator}
  \Big\langle B_x B_y \prod_{i=1}^4 \mathscr V_i \Big\rangle
    = (- 2 \I \zeta) (-1) \, \big\langle c \bar c V_1 \, c \bar c V_2 \, c \bar c V_3 \, V_4(z_4) \big\rangle
    = 2 \I \zeta \, \big\langle c \bar c V_1 \, c \bar c V_2 \, c \bar c V_3 \, V_4(z_4) \big\rangle.
\end{equation}

\paragraph{Step 5: integration.}
With the prefactor $\kappa / 2\pi$, the positive orientation $\int \dd x \wedge \dd y = \int \dd x \, \dd y$ and $\dd^2 z = 2 \, \dd x \, \dd y$, the four-point form integrates to
\begin{equation}
  \label{eq:conv:measure-N4}
  \int_{\mc S_{0,4}} \omega_2 \Big\vert_{\text{on shell}}
    = \frac{\kappa}{2\pi} \, 2 \I \zeta \int \dd x \, \dd y \, \big\langle \cdots V_4(z_4) \big\rangle
    = \frac{\kappa \, \I \zeta}{2\pi} \int \dd^2 z_4 \, \big\langle \cdots V_4(z_4) \big\rangle
    = \frac{\kappa \, \I \zeta}{2\pi} \, A^{\rm FP}_{0,4},
\end{equation}
where $A^{\rm FP}_{0,4} := \int \dd^2 z_4 \, \langle \prod_{i \le 3} c \bar c V_i \, V_4 \rangle$ is the on-shell four-point amplitude with three punctures fixed by the Faddeev--Popov procedure.

\paragraph{The observable sign.}
Only the product
\begin{equation}
  \label{eq:conv:xi}
  \xi
    := \kappa \, \I \zeta
    \in \{ +1, -1 \}
\end{equation}
enters the amplitude.
The four combinations of a prefactor and a contour convention give:
\begin{center}
  \begin{tabular}{lcccc}
    \toprule
    Prefactor & Contour reading & $\kappa$ & $\zeta$ & $\xi$ \\
    \midrule
    $\I / 2\pi$~\cite{Belopolsky:1995:OffshellClosedString} & literal, \eqref{eq:conv:zeta-literal} & $+ \I$ & $-1$ & $+1$ \\
    $\I / 2\pi$~\cite{Belopolsky:1995:OffshellClosedString} & $+ \Res_{\bar w}$ & $+ \I$ & $+1$ & $-1$ \\
    $(2\pi \I)^{-1}$~\cite{Erbin:2021:StringFieldTheory} & literal, \eqref{eq:conv:zeta-literal} & $- \I$ & $-1$ & $-1$ \\
    $(2\pi \I)^{-1}$~\cite{Erbin:2021:StringFieldTheory} & $+ \Res_{\bar w}$ & $- \I$ & $+1$ & $+1$ \\
    \bottomrule
  \end{tabular}
\end{center}
The four rows are not four claims about one quantity: they are two claims, $\xi = \pm 1$, reached by two independent binary choices.
The factor $(-1)^{N-3}$ between the printed prefactors is reproduced by flipping either $\kappa$ or $\zeta$, so comparing the two prefactors without comparing the two contour conventions is not a well-posed comparison.
The prefactor of~\cite{Erbin:2021:StringFieldTheory} read with its own anticlockwise rule (third row) and the prefactor of Belopolsky and Zwiebach read with the residue convention (second row), which is the reading used with their prefactor in the explicit evaluation of the quartic vertex with $\group{SL}(2, \C)$ local coordinates, give the same $\xi = -1$.

\paragraph{Any number of points.}
For $N$ punctures the moduli are the positions $z_4, \ldots, z_N$.
In the frame $(x_4, y_4, x_5, y_5, \ldots)$, the top form has the component
\begin{equation}
  \omega_{2(N-3)}(\pd_{x_4}, \pd_{y_4}, \ldots)
    = \Big( \frac{\kappa}{2\pi} \Big)^{N-3}
      \Big\langle (B_{x_4} B_{y_4}) (B_{x_5} B_{y_5}) \cdots \prod_i \mathscr V_i \Big\rangle.
\end{equation}
Each pair $B_{x_j} B_{y_j}$ is a pair of contours around the puncture $j$ only and is Grassmann even, so it commutes with every other pair and every $\mathscr V_i$, and it contributes the factor $2 \I \zeta$ of~\eqref{eq:conv:four-point-correlator} on its own puncture.
Hence
\begin{equation}
  \label{eq:conv:measure-N}
  \int_{\mc S_{0,N}} \omega_{2(N-3)} \Big\vert_{\text{on shell}}
    = \Big( \frac{\xi}{2\pi} \Big)^{N-3} A^{\rm FP}_{0,N},
  \qquad
  A^{\rm FP}_{0,N}
    = \int \prod_{j=4}^N \dd^2 z_j \,
      \Big\langle \prod_{i=1}^3 c \bar c V_i \prod_{j=4}^N V_j \Big\rangle.
\end{equation}
The sign is a pure power of $\xi$, with no further $N$-dependent sign from the reordering of the antighosts.
Three checks support the formula.
At $N = 3$ there is no modulus and no antighost insertion, the prefactor is one, and both transcriptions agree on the cubic amplitude, as they must.
The amplitude and the integrated form are proportional by a real factor, since no convention introduces a stray phase, so $\xi$ must be real and $\kappa \in \{ \pm \I \}$ for $\zeta = \pm 1$: this is why both sources carry a factor $\I$, but it does not choose the sign.
The prefactor is dimensionless in both transcriptions, and the dispute is a pure phase.

\paragraph{The value of $\xi$.}
The computation fixes $\int \omega$ in terms of $A^{\rm FP}$ up to the sign $\xi$ per modulus, and it fixes the magnitude $1 / 2\pi$ and the reality of $\xi$.
It does not fix the sign: that requires the sign of the normalization of the string field theory amplitude relative to the Faddeev--Popov one, which is fixed by comparing the Feynman region of the four-point amplitude with the exchange term of the action, or by the residue of the four-tachyon amplitude at the tachyon pole.
That comparison, made with the explicit quartic vertex, gives $\xi = -1$; we take it as an input here and do not repeat it.
Everything else in this section is independent of it: the identity~\eqref{eq:conv:measure-N}, the reading of the table, the value $\zeta = -1$ of the anticlockwise rule, and the agreement of the two transcriptions.

\paragraph{Consequences for the two references and for this write-up.}
The two printed prefactors are not statements about the same quantity: a prefactor without a declared antiholomorphic contour convention determines no physics.
Read with the contour rule it states, the measure of~\cite{Erbin:2021:StringFieldTheory} gives $\xi = -1$, the same value as the prefactor of Belopolsky and Zwiebach read with the residue convention, so the two references agree and neither needs a correction; what the first lacks is an explicit statement, next to the measure, that $\dd \bar z$ is the differential along the anticlockwise contour, so that $\oint_C \frac{\dd \bar z}{2\pi \I} \, \bar z^{-n-2} = - \delta_{n,-1}$.
Without it, a reader who adopts the residue convention obtains the wrong sign per complex modulus.
For this write-up, with $\xi = -1$, the constant of the quartic contact term is negative on the potential side, and the sign left open in the Schwinger representation of the inverse of the quadratic operator of $\Sigma_2$ (\cref{ch:worldsheet-dictionary}) is fixed: the representation without a sign condition applies.
The identity~\eqref{eq:conv:measure-N} is checked by the companion script on the measure orientation: the contour integral~\eqref{eq:conv:zeta-literal} by numerical quadrature, the steps~\eqref{eq:conv:Bxy}--\eqref{eq:conv:ghost-strip} in explicit finite-dimensional representations of the ghost modes, and the pure power at $N = 5$; the value of $\xi$ is not checked there.


\subsection{The symmetry type of the strands}
\label{sec:conv:strands}

\status{Verified in a finite model}\record{076-hwedge-convention}

The paper states, in its convention appendix, that the components of a tensor in $\Hn{n}$ are graded symmetric in the parity of the strands for the arguments of vertices and products, and in the shifted parity for the strands that the $n$-string pairing sets against a field.
This section derives that rule and shows that under it every pairing and product of the theory is well defined and nonzero where it must be, while the physical parity on every strand makes the multiplier term vanish.

\paragraph{Input strands and output strands.}
The one-string pairing is graded symmetric on its support, and the $n$-string pairing is its tensor power with Koszul signs,~\eqref{p:eq:conv:pairing-n}; the parity is $\varepsilon = \Ngh \bmod 2$.
A strand is an \emph{input} when it is an argument of a vertex or a product: the strands of $\Psi^n$, $\Sigma_n$, the gauge parameters, the equations of motion.
It is an \emph{output} when the $n$-string pairing sets it against an input strand: the outputs of the products $\sprod{m}{n}$ and the strands of the multipliers $\Theta_k$, which multiply $\Sigma_k - \Psi^k$ through the $k$-string pairing.
The argument is finite linear algebra and needs no hypothesis.

\paragraph{Input strands.}
The literature fixes only the input strands.
In~\cite{Erler:2014:NSNSSectorClosed} the degree of a closed string state is its Grassmann parity and the wedge product obeys $\Phi_1 \wedge \Phi_2 = (-1)^{\deg \Phi_1 \deg \Phi_2} \, \Phi_2 \wedge \Phi_1$; there, $\Hn{n}$ is only the domain of the multilinear products, and no map has several outputs.
Münster and Sachs~\cite{Muenster:2012:HomotopyAlgebrasString-kn} build the graded symmetric algebra of the closed string field shifted by two degrees, an even shift, so it is again graded symmetric in the physical parity on every strand.
On inputs the rule is therefore theirs: even strands have symmetric components, odd strands antisymmetric ones.
Neither reference pairs a multi-string tensor against another, so neither assigns a symmetry type to output strands; that type is fixed by the pairing, as follows.

\paragraph{Output strands are graded symmetric in the shifted parity.}
Raising a leg of parity $\varepsilon_b$ of a vertex through the pairing gives an output of parity $\varepsilon_a = \varepsilon_b + 1$, since the pairing joins states of opposite parity.
Exchanging two outputs $a_i \leftrightarrow a_j$ of $\sprod{m}{n}$ is exchanging the two corresponding legs $b_i \leftrightarrow b_j$ of the vertex, which costs $(-1)^{\varepsilon_{b_i} \varepsilon_{b_j}} = (-1)^{(\varepsilon_{a_i} + 1) (\varepsilon_{a_j} + 1)}$.
The outputs are thus graded symmetric in the shifted parity, for inputs of any parity.
At ghost number three every output strand is odd and the shifted parity is even: the outputs have plainly symmetric components.

\paragraph{Only the symmetric part pairs with $\Sigma_k$.}
Take $\Sigma \in \Hn{k}$ with strands of ghost number two and a tensor $\Theta$ with strands of ghost number three.
Since the strands of $\Sigma$ are even, the Koszul sign of~\eqref{p:eq:conv:pairing-n} is one, and in components
\begin{equation}
  \label{eq:conv:strands-contraction}
  \pair[k]{\Theta}{\Sigma}
    = \sum \Theta^{b_1 \cdots b_k} \, \prod_{i=1}^k \pair{e_{b_i}}{e_{a_i}} \, \Sigma^{a_1 \cdots a_k}.
\end{equation}
The product of pairings is invariant when the same permutation acts on the $a$ and the $b$ indices.
If $\Theta$ had antisymmetric components, which is the graded symmetry in the physical parity on odd strands, a transposition of two strands would change the sign of the sum while leaving $\Sigma^{a_1 \cdots a_k}$ unchanged: the sum would vanish identically, whatever the signs of the one-string pairing.
The smallest instance is $k = 2$, with $\Theta = B_1 \otimes B_2 - B_2 \otimes B_1$ and $\Sigma = A_1 \otimes A_2 + A_2 \otimes A_1$:
\begin{equation}
  \pair[2]{\Theta}{\Sigma}
    = \pair{B_1}{A_1} \pair{B_2}{A_2}
      + \pair{B_1}{A_2} \pair{B_2}{A_1}
      - \pair{B_2}{A_1} \pair{B_1}{A_2}
      - \pair{B_2}{A_2} \pair{B_1}{A_1}
    = 0.
\end{equation}
So the part of a ghost-number-three tensor that pairs with $\Sigma_k$ is the part with plainly symmetric components, which is the image of the raising above.
With the physical parity on every strand, the multiplier term $\pair[k]{\Theta_k}{\Sigma_k - \Psi^k}$ would vanish identically and the cubic action would contain no constraint.
The pairing thus forces the shifted parity on output strands.

\paragraph{Two odd inputs.}
The product $\ell_{n+2}$ is graded symmetric in the parity of its inputs.
On two gauge parameters, both odd, it is antisymmetric,
\begin{equation}
  \ell_{n+2}(\Lambda_1, \Lambda_2, \Psi^n)
    = - \ell_{n+2}(\Lambda_2, \Lambda_1, \Psi^n),
\end{equation}
generically nonzero, and zero at $\Lambda_1 = \Lambda_2$, as the commutator $[\delta_\Lambda, \delta_\Lambda] = 0$ of a transformation with itself requires.
This is the rule on input strands of~\cite{Erler:2014:NSNSSectorClosed}.

\paragraph{Mixed strands.}
The two gradings give the exchange signs $(-1)^{\varepsilon_i \varepsilon_j}$ and
\begin{equation}
  (-1)^{(\varepsilon_i + 1) (\varepsilon_j + 1)}
    = (-1)^{\varepsilon_i \varepsilon_j + \varepsilon_i + \varepsilon_j + 1},
\end{equation}
which are equal if and only if $\varepsilon_i + \varepsilon_j = 1$.
A tensor whose two strands have opposite parities has the same symmetry type in both gradings, and its role need not be specified.
On two strands of the same parity the two types are different subspaces, of dimensions $d (d + 1) / 2$ and $d (d - 1) / 2$ for a $d$-dimensional space of strands, and no re-signing $T^{ab} \mapsto (-1)^{f(a,b)} \, T^{ab}$ maps one onto the other, since on the diagonal $a = b$ it would need $(-1)^0 = -1$.
The distinction between input and output strands is therefore a property of the arrays, which no choice of signs removes.

\paragraph{The rule.}
Components are graded symmetric in the physical parity on input strands, as in~\cite{Erler:2014:NSNSSectorClosed,Muenster:2012:HomotopyAlgebrasString-kn}, and in the shifted parity on output strands, which are the dual of the input strands under the multi-string pairing and the image of the raising of a vertex leg.
At ghost number three the output strands have symmetric components; for mixed tensors the role is irrelevant; the outputs of the higher products of the IBL structure (\cref{ch:ibl}) are output strands of any ghost number and follow the same rule.
One symbol, $\Hn{n}$, serves for both roles.
The physical content is that a constraint imposed through the multi-string pairing is nontrivial only if the multiplier has the symmetry dual to that of the field it constrains.
The statements on the raised outputs, on the vanishing~\eqref{eq:conv:strands-contraction} for antisymmetric multipliers, on the product of two gauge parameters and on the signs and dimensions of the last paragraph are checked in a finite model with three states at ghost number two and three at ghost number three, a random one-string pairing with the symmetry above and a random graded-symmetric quartic vertex; the vanishing check gives a nonzero pairing for a multiplier with symmetric components, so it can fail.
The reading of the two references is not machine-checked.
